\section{Exact colimits and compact inverse diagrams}
\label{sec:C}

We use $\mathbb T=\R/\Z$ with its usual compact topology. If $M$ is a discrete abelian group, $M^\vee=\Hom_{\Z}(M,\mathbb T)$ has the topology of pointwise convergence. If $K$ is compact Hausdorff abelian, $K^\wedge=\Hom_{\mathrm{cts}}(K,\mathbb T)$ is discrete. In all occurrences of $\lim^i$ below, the derived functor is taken in the category of diagrams of \emph{abstract} abelian groups; continuity is a requirement on the original diagram, not an alternative definition of its derived limit.

\subsection{Duality and torsors}

\begin{theorem}[C.1; Proved: the bar--cobar duality]\label{C:duality}
For every small category $A$, every $M:A\to\Ab$, and every integer $i\geq0$, there is a natural isomorphism of abstract abelian groups
\[
 \lim_{A^{\op}}^{i}M^\vee\ \cong\ \Hom_{\Z}(L_i\colim_A M,\mathbb T).
\]
For $M:A\to\F_p\text{-}\mathrm{Mod}$ the analogous duality is
\[
 \lim_{A^{\op}}^{i}\Hom_{\F_p}(M,\F_p)
 \cong\Hom_{\F_p}(L_i\colim_A M,\F_p),
\]
where the dual groups carry their profinite vector-space topologies. For a $\Q$-module diagram, the first displayed formula holds with $L_i\colim_A$ computed in $\Q$-modules and with the Pontryagin dual of its underlying additive group. These statements include empty and disconnected $A$.
\end{theorem}
\begin{proof}
Write the unnormalized simplicial replacement as
\[
 C_n(A;M)=\bigoplus_{a_0\xrightarrow{f_1}\cdots\xrightarrow{f_n}a_n}M(a_0).
\]
Its zeroth face applies $M(f_1)$, its intermediate faces compose adjacent arrows, and its last face deletes the last arrow. The alternating face map is the differential. This complex computes $L_*\colim_A M$. One justification that does not assume the groups $M(a)$ are projective is to use the simplicial resolution by the evaluation--induction adjunction: an induced diagram $\Z[A(a,-)]\otimes M(a)$ is colimit-acyclic, because its standard replacement is the chain complex of the contractible category $a/A$ with the specified constant coefficient group. The resulting acyclic resolution computes the derived functor. Equivalently, resolve every $M(a)$ by free groups and compare the two filtrations of the double complex.

Applying $\Hom_{\Z}(-,\mathbb T)$ turns this direct-sum complex into
\[
 C^n(A^{\op};M^\vee)=\prod_{a_0\to\cdots\to a_n}M(a_0)^\vee.
\]
This is the Roos cobar complex, with precisely its usual transition maps. In degree zero its differential is
\[
 (\partial t)_f=M(f)^\vee(t_b)-t_a\qquad(f:a\to b).
\]
Since $\mathbb T$ is injective as an abelian group, dualization is exact and identifies cohomology with the dual of homology. This proves the first formula, including naturality. Over $\F_p$ the same argument uses the exact vector-space dual and the product topology. Every profinite $\F_p$-vector space is the continuous dual of a discrete vector space. For $\Q$, the integral and rational bar complexes of a rational diagram coincide after forgetting scalars. Its compact dual is generally a solenoidal compact group, not a vector-space dual with the discrete topology. No comparison with such a discrete dual is being asserted.
\end{proof}

\begin{theorem}[C.2; Proved: exactness is the compact torsor property]\label{C:torsors}
For every small category $A$, the following conditions are equivalent:
\begin{enumerate}[label=(\roman*)]
\item $\colim_A:\Ab^A\to\Ab$ is exact;
\item $\lim_{A^{\op}}$ takes objectwise surjective morphisms of compact Hausdorff abelian diagrams to surjections;
\item $\lim^1_{A^{\op}}K=0$ for every compact Hausdorff abelian diagram $K$;
\item every compact torsor diagram on $A^{\op}$ has a point in its limit, namely (CT).
\end{enumerate}
Replacing compact abelian groups by profinite $\F_p$-vector spaces gives exactness over $\F_p$ and $(\mathrm{CT})_p$. Replacing them by Pontryagin duals of discrete $\Q$-vector spaces gives exactness over $\Q$ and $(\mathrm{CT})_{\Q}$.
\end{theorem}
\begin{proof}
Pontryagin duality is an exact contravariant equivalence between discrete abelian groups and compact Hausdorff abelian groups; in particular, a monomorphism of discrete groups dualizes to a continuous surjection. The coend description, or degree zero in Theorem~\ref{C:duality}, identifies the dual of a colimit with the limit of the dual diagram. Thus (i) and (ii) are equivalent. Theorem~\ref{C:duality}, the fact that $\mathbb T$ is a cogenerator, and the criterion $L_1\colim=0$ for exactness give (i)$\Leftrightarrow$(iii).

For the torsor assertion, choose a point $t_a\in T(a)$ for each object of a torsor diagram under $K$. The simply transitive compact action identifies $T(a)$ homeomorphically with $K(a)$ by $x\mapsto x+t_a$: the orbit map is a continuous bijection from a compact space to a Hausdorff space. For $f:a\to b$ there is a unique
\[
 c_f=T(f)t_b-t_a\in K(a).
\]
For $a\xrightarrow f b\xrightarrow g c$, equivariance gives
\[
 c_{gf}=c_f+K(f)c_g.
\]
Consequently $(c_f)$ is a degree-one cocycle in the above cobar complex. Replacing $t_a$ by $t_a+x_a$ changes $c$ to $c+\partial x$. A compatible family is exactly a choice for which this new cocycle is zero.

Conversely, any degree-one cocycle $c_f\in K(a)$ defines a compact torsor diagram with $T(a)=K(a)$ as a space and
\[
 T(f)(x)=K(f)x+c_f\quad(x\in K(b)).
\]
The cocycle equation is precisely the functoriality equation, and each map is continuous. Thus every class in $\lim^1K$ occurs, and it is zero exactly when its torsor has a section. This proves (iii)$\Leftrightarrow$(iv). The same proof works in the two specified dual coefficient categories, which are closed under kernels. For example, a point in the limit of a quotient diagram has a torsor of lifts under its kernel, explaining (ii) directly in terms of (iv).
\end{proof}

\subsection{Finite coefficient tests and their limitations}

\begin{proposition}[C.3; Proved and Refuted: finite vector torsors]\label{C:finite-tests}
Fix a prime $p$. If $A$ has finite hom-sets, exactness of $\colim_A$ over $\F_p$ is equivalent to the assertion that every torsor diagram under \emph{finite} $\F_p$-vector spaces has nonempty limit. The forward implication holds for every small category. The converse is false without the hom-set hypothesis, even for a nonempty connected category satisfying condition I.
\end{proposition}
\begin{proof}
The forward implication is Theorem~\ref{C:torsors}. Conversely, the finite-torsor assertion and its cocycle proof say that $T_1(M)=0$ for every $M:A\to\F_p\text{-}\mathrm{Mod}$ that is finite dimensional at each object. Every vector-space diagram is the filtered union of its subfunctors generated by finitely many elements at finitely many objects. A subfunctor generated by $m_j\in M(a_j)$, $1\leq j\leq s$, is at $b$ spanned by
\[
 \{M(f)m_j:1\leq j\leq s,\ f\in A(a_j,b)\}.
\]
This set is finite under the hom-set assumption. The bar complex and its homology commute with filtered colimits, so $T_1$ vanishes on every diagram. Notice that no \emph{uniform} bound on these dimensions has been used.

For the counterexample take the one-object category $B\Q$ of the additive discrete group $\Q$. A homomorphism from $\Q$ to a finite group is trivial: its image is abelian and divisible, whereas a finite divisible abelian group is zero. Thus every action on a finite set is trivial. A finite torsor diagram on $(B\Q)^{\op}$ is an action of $\Q$ by permutations of a finite nonempty set, so every point is compatible.

Let $k=\F_p$, let $P=k[\Q]$ be the left regular module, and define the natural left-module endomorphism
\[
 D:P\longrightarrow P,\qquad [q]\longmapsto[q+1]-[q].
\]
It is injective: on each coset of $\Z\subset\Q$, a finitely supported translation-invariant coefficient function is zero. Its map on coinvariants is zero, and $P$ is representable. Hence $T_1(\coker D)=k\ne0$. The same construction works over every nonzero unital coefficient ring. Condition I holds because the category is a group. This proves the failure of the proposed unrestricted converse.
\end{proof}

\begin{proposition}[C.4; Refuted: uniformly bounded finite sets characterize quasi-filteredness]\label{C:bounded-tests}
For a nonempty connected category satisfying I, the condition that every contravariant diagram of nonempty finite sets of uniformly bounded cardinality has nonempty limit does not imply quasi-filteredness. It does not imply an integral affine equalizer even for a single pair.
\end{proposition}
\begin{proof}
Again take $A=B\Q$. Every finite-set action is trivial, so its limit is the same nonempty finite set. For the parallel endomorphisms $0,1$ of its unique object, an affine equalizer would be a finitely supported $t\in\Z[\Q]$ with $t(1-0)=0$ and augmentation one. The translation argument in Proposition~\ref{C:finite-tests} forces $t=0$, a contradiction. The same proof shows that this pair has no affine equalizer over any nonzero ring. The category is small, nonempty, connected, and satisfies I, exactly as requested.
\end{proof}

\begin{theorem}[C.5; Proved: the profinite test]\label{C:profinite}
For every small category $A$, the following are equivalent:
\begin{enumerate}[label=(\roman*)]
\item $\lim^1_{A^{\op}}K=0$ for every profinite abelian diagram $K$;
\item $\colim_A$ is exact over $\F_p$ for every prime $p$;
\item every objectwise torsion diagram $M:A\to\Ab$ has $T_i(M)=0$ for every $i>0$;
\item $\lim^i_{A^{\op}}K=0$ for every profinite abelian diagram and every $i>0$.
\end{enumerate}
\end{theorem}
\begin{proof}
The Pontryagin duals of profinite abelian groups are exactly the discrete torsion groups. This follows directly in one direction because a continuous character on a profinite group has finite image, and in the other by writing a torsion group as the union of its finite subgroups. Theorem~\ref{C:duality} therefore proves (i)$\Rightarrow$(ii) and (iii)$\Rightarrow$(iv)$\Rightarrow$(i).

Assume (ii). The integral bar complex of an $\F_p$-module diagram is its field bar complex, so all its positive homology groups vanish. For a $p$-primary torsion diagram, filter by the subfunctors $M[p^n]$. Their successive quotients are killed by $p$. Closure of acyclic diagrams under extensions, followed by filtered unions, proves acyclicity of the $p$-primary part. The full torsion diagram is the direct sum of these parts. Direct sums, filtered colimits, and the bar differential commute, and these colimits are exact in abelian groups. This proves (iii).
\end{proof}

\subsection{Arithmetic of affine equalizers}

For a finite family $a_0,\ldots,a_k:p\to q$, set
\[
 P=\Z[A(q,-)],\qquad E=\Z[A(p,-)]^k,\qquad
 D(t)=(ta_i-ta_0)_{1\leq i\leq k}.
\]
When $A$ has common successors, the union of augmentation images of $\ker D(r)$ is an ideal $I\subseteq\Z$: to add elements at different targets, move them to a common successor. We use $\ker D$ as a subfunctor of $P$, not merely as a collection of unrelated kernels.

\begin{proposition}[C.6; Proved: pair equalizers and exceptional primes]\label{C:arithmetic-pairs}
Suppose $A$ has common successors.
\begin{enumerate}[label=(\alph*)]
\item For a pair $a_0,a_1$, existence of an affine equalizer over one $\F_p$ at a target $r$ implies existence of a rational affine equalizer at that same target. If the pair has an affine equalizer over every $\F_p$, possibly at different targets, it has an integral affine equalizer.
\item For any finite family and a target $r$ with $A(q,r)$ and $A(p,r)$ finite, if a prime-field affine equalizer exists at $r$ but no rational one does, the prime divides the torsion subgroup of $\coker D(r)$.
\item If $A$ has a weakly terminal object and is exact over every $\F_p$, it is filtered and hence exact over $\Z$. The same conclusion holds if $A$ is total and has a weakly terminal object. These conclusions do not require finite hom-sets; in particular every total one-object category is filtered.
\end{enumerate}
\end{proposition}
\begin{proof}
For (a), $D(r)$ is the integral incidence map of the directed graph whose vertices are $A(p,r)$ and whose edges are $g\in A(q,r)$, directed from $ga_0$ to $ga_1$. Its cokernel is the free abelian group on the connected components. This remains true for an infinite graph: every relation has finite support and the quotient identifies precisely the vertices in the same finite zigzag component. Thus the cokernel is torsion-free, and reduction of the kernel modulo $p$ surjects onto the kernel of the reduced incidence map. A field equalizer lifts to an integral kernel element of augmentation congruent to one modulo $p$, which is nonzero as an integer. Dividing by its augmentation proves rational solvability at the same target. If this works for every prime, $I$ is divisible by no prime and is nonzero; hence $I=\Z$.

For (b), at a prime not dividing the finite torsion subgroup of $\coker D(r)$ the same kernel-lifting argument applies. Its lift has nonzero augmentation and therefore gives a rational affine equalizer.

For (c), exactness over any nonzero field implies condition I by the coefficient-extraction proof of \citet[Lemma 3.2 and the coefficient-ring remark after Proposition 3.3]{IsbellMitchell1976}. Explicitly, an affine commutative square has equal sums of ordinary composites with augmentation one on both sides, so a nonzero term on one side must also occur on the other, giving an ordinary square. Part (a) gives condition IV. Thus $A$ is quasi-filtered. \citet[Corollary 4.4]{IsbellMitchell1976} states and proves that a connected category with I, IV and a weak terminal object is filtered. We have opened and checked that proof: IV filters every finite family of endomorphisms (Corollary 4.3); moving a pair to the weak terminal object and applying I reduces its ordinary equalization to that endomorphism assertion. This also proves the stronger conclusion claimed here. A total category is quasi-filtered, so the same opened Corollary4.4 proves the assertion about total categories and, in particular, one-object categories.

Here is an independent direct proof of (c). First let $M$ be a monoid
exact over every $\F_p$. For $x,a\in M$, an affine equalizer gives
$t\in\F_p[M]$ with augmentation one and $txa=tx$. Put $z=tx$.
Every point of the finite support of $z$ has a predecessor in that
support under right multiplication by $a$. Following predecessors
shows that every support point lies on a finite cycle; its coefficient
is constant around the cycle. Since the augmentation is one, at least
one such cycle has length $r$ prime to $p$. Thus some left multiple
$yx$ satisfies $yx a^r=yx$. Start with $p=2$. If the obtained period
$r$ exceeds one, choose a prime $p\mid r$ and apply the same argument
to the new $x'=yx$. The resulting left multiple retains period $r$
and acquires a period $s$ prime to $p$, hence has period
$\gcd(r,s)<r$. Repetition gives a left multiple with period one.
We have proved that for every $x,a$ there is $y$ with $yxa=yx$.
Applying this first to $x=1,a$ and then to the resulting $x=y,b$
equalizes any two $a,b\in M$, so $M$ is filtered.

Now let $c$ be weakly terminal in $A$. Postcomposing affine equalizers
with arrows to $c$ shows that $\operatorname{End}_A(c)$ is exact over
each $\F_p$, so it is filtered by the preceding argument. Given
$f,g:p\to q$, choose $h:q\to c$. A prime-field affine equalizer,
again followed by an arrow to $c$, is a finite affine combination
$t$ of endomorphisms of $c$ with $thf=thg$. Filteredness of the
endomorphism monoid supplies one endomorphism $k$ such that all
elements in the support of $kt$ are the same arrow $u$.
Its augmentation is one, so $kt=[u]$ and $uhf=uhg$ as ordinary
arrows. This proves filteredness without an external reduction.
\end{proof}

\begin{proposition}[C.7; Proved: the rational-only obstruction]\label{C:rational-gate}
Assume $A$ has common successors and is exact over $\F_p$ for every prime. Then it is quasi-filtered. For every objectwise torsion-free $M:A\to\Ab$, each $T_i(M)$, $i>0$, is uniquely divisible. For every finite parallel family its augmentation ideal satisfies $I=0$ or $I=\Z$. Consequently rational exactness and integral exactness are equivalent under the all-prime hypothesis.

If a family has $I=0$, define the subfunctor
\[
 S(r)=\{s\in\Q\otimes P(r):D(r)s\in E(r)\}.
\]
Then augmentation gives isomorphisms
\[
 \colim_A S\cong\Q,\qquad
 \colim_A (\coker D)_{\mathrm{tor}}\cong\Q/\Z.
\]
If hom-sets are finite, the torsion cokernel at each individual target has only finitely many prime divisors. Thus the possible failure is necessarily distributed among unboundedly many targets.
\end{proposition}
\begin{proof}
Condition I and IV were proved above. For an objectwise torsion-free diagram the sequence
\[
 0\longrightarrow M\xrightarrow{p}M\longrightarrow M/pM\longrightarrow0
\]
is exact. All positive derived colimits of its last term vanish by Theorem~\ref{C:profinite}. The long exact sequence says multiplication by $p$ is bijective on $T_i(M)$ for every $i>0$. This holds for every prime, proving unique divisibility.

Put $J=\im D$ and $C=\coker D$. Right exactness applied to $\ker D\to P\to J$ gives $\colim J=\Z/I$, because the colimit of a representable is $\Z$. The induced map $\colim J\to\colim E=\Z^k$ is zero. As $E$ is projective, the long exact sequence gives
\[
 T_1(C)\cong\Z/I.
\]
The torsion subfunctor $C_{\mathrm{tor}}$ is acyclic. The sequence $0\to C_{\mathrm{tor}}\to C\to C/C_{\mathrm{tor}}\to0$ therefore injects $T_1(C)$ into the rational vector space $T_1(C/C_{\mathrm{tor}})$. A nonzero finite cyclic group cannot embed in a rational vector space. Thus $I=0$ or $\Z$. If rational exactness holds, each finite family has a rational kernel element of augmentation one, which after clearing denominators makes $I\ne0$; therefore $I=\Z$. The affine criterion proves integral exactness. Its converse is immediate by restriction of scalars.

We prove the saturation assertions without presuming exactness on non-torsion diagrams. The quotient $S/P$ is torsion and is a subfunctor of $(\Q\otimes P)/P$. Exactness on torsion diagrams makes
\[
 \colim(S/P)\longrightarrow\colim((\Q\otimes P)/P)=\Q/\Z
\]
injective. Also $T_1(S/P)=0$, so $\Z=\colim P\to\colim S$ is injective. Comparing the two short exact sequences shows that augmentation identifies $\colim S$ with a subgroup of $\Q$ containing $\Z$.

Exactness on torsion diagrams implies exactness over $\Z/p^m\Z$. Its affine criterion supplies $t\in P(r)$ with $D(t)\in p^mE(r)$ and $\aug t\equiv1\pmod {p^m}$. Then $t/p^m\in S(r)$. Its augmentation, together with the integers, generates $1/p^m$ in $\colim S\subset\Q$. Varying $p,m$ proves $\colim S=\Q$ and $\colim(S/P)=\Q/\Z$.

Finally write $K_{\Q}=\ker(\Q\otimes D)$ and $K=\ker D$. There is an exact sequence of torsion diagrams
\[
 0\longrightarrow K_{\Q}/K\longrightarrow S/P
   \longrightarrow C_{\mathrm{tor}}\longrightarrow0.
\]
The last map sends the class of $s$ to the class of $Ds$; its kernel is exactly $(P+K_{\Q})/P=K_{\Q}/K$. Because $I=0$, every element of $K_{\Q}$ has augmentation zero. The induced map from $\colim(K_{\Q}/K)$ into $\colim(S/P)\subset\Q/\Z$ is both injective (torsion exactness) and zero. Thus the last colimit is $\Q/\Z$. The finite-target assertion is the elementary finite-matrix torsion bound.
\end{proof}

\subsection{An exceptional-prime finite pattern}
\begin{proposition}[Proved: modular affine solvability need not lift at one target]
\label{C4:exception}
There is a finite category with three objects, finite hom-sets and common
successors, and a triple of parallel arrows with an affine equalizer over
$\mathbb F_2$ but none over $\mathbb Q$. This assertion concerns the
specified family, not exactness of the category over $\mathbb F_2$.
\end{proposition}
\begin{proof}
Let the objects be $p,q,r$. In addition to identities, let the arrows
be $a_0,a_1,a_2:p\to q$, $b_1,\ldots,b_7:q\to r$, and
$c_0,\ldots,c_5:p\to r$. Define $b_j a_i=c_{e_{j,i}}$, where the
ordered triples $e_j$ are
\[
 012,\quad013,\quad015,\quad025,\quad124,\quad135,\quad234.
\]
These are all nonidentity compositions; identities provide the rest,
so associativity is immediate and is also checked exhaustively by the
program. Every object maps to $r$.

Over $\mathbb F_2$, take
$t=b_1+b_3+b_5+b_6+b_7$, of augmentation one. Its three marginals
coincide modulo two, as direct substitution shows.
For a rational certificate in the other direction, assign values on
$c_0,\ldots,c_5$ to two functionals:
\[
 \lambda_1=(-3/2,-1/2,-1/2,0,0,0),\qquad
 \lambda_2=(0,-1/2,0,0,1/2,0).
\]
Every listed triple $e_j=(x,y,z)$ satisfies
$\lambda_1(y)-\lambda_1(x)+\lambda_2(z)-\lambda_2(x)=1$.
Thus for $D(t)=(ta_1-ta_0,ta_2-ta_0)$,
$(\lambda_1,\lambda_2)D=\operatorname{aug}$ over $\mathbb Q$.
A rational kernel element has augmentation zero, so there is no
affine equalizer at $r$. At $q$ the only possible affine endomorphism
is the identity, which does not equalize the arrows; there are no
arrows $q\to p$. This proves the assertion at all targets.

The code \texttt{C4\_exception\_certificate.py} verifies the modular
vector, the rational functional, and every composable triple in this
19-arrow category. It also records why this is not a counterexample
to the central all-prime question: two distinct $b_j$ cannot be
equalized over any ring, since $r$ has only its identity endomorphism.
\end{proof}

\subsection{Contractible Kan simplex categories are Ramsey}
\begin{theorem}[Proved: the contractible Kan candidate is already exact]
\label{C4:contractible-kan-ramsey}
Let $K$ be a nonempty contractible Kan complex. Then the category
$\Delta_{\mathrm{inj}}/K$ has the morphism Ramsey property. Consequently
its colimit functor is exact over every unital ring. In particular,
choosing a contractible Dold--Kan complex from an acyclic but nonsplit
chain complex does not give a total nonexact example in this family.
\end{theorem}
\begin{proof}
Fix objects $A,B$ and a finite number $c$ of colours. If there is no
arrow $A\to B$, the conclusion is immediate with target $B$. Otherwise
let $L$ be the finite relational language with one relation $R_\sigma$
for each simplex label occurring on a nonempty face of $A$ or $B$;
its arity is the number of vertices of that face. Equal labels give
one symbol. This set of symbols is closed under all face maps.

Consider finite $L$-structures with the following rules. Every vertex
has exactly one vertex label. Every relation tuple has distinct
entries; it has at most one label of its arity. If an ordered tuple
has label $\sigma$, every ordered subtuple has the corresponding
face label of $\sigma$. No global order on these tuples is required
at this stage, and different orderings of the same vertex set may
be labelled independently, subject only to their common ordered
subtuples. This is a hereditary free amalgamation class: in an
amalgam all relation tuples come from one factor, and their subtuples
come from that factor as well; labels on a tuple in both factors
already agree on the common induced substructure. There are no new
cross-factor relation tuples.

Add an arbitrary linear order. Theorem~1.3 of
Evans--Hubi\v cka--Ne\v set\v ril \cite{D-EHN2021} gives a finite
ordered member $C$ with $C\longrightarrow(B)^A_c$, where $A$ and $B$
are encoded by labelling precisely their increasing tuples. Remove
from $C$ every relation tuple that is not increasing in its added
linear order. The remaining labels give a downward closed simplicial
subcomplex $L_C$ of the full simplex on $V(C)$ and a coherent
simplicial map
\[
 L_C\longrightarrow K.
\]
Deleting the nonincreasing tuples does not change any existing
induced copy of $A$ or $B$, since these structures have no such
tuples. In particular the specific copies used in the Ramsey
conclusion remain available.

The map $K\to *$ is a trivial Kan fibration: it is a Kan fibration
and a weak equivalence. Thus it has the right lifting property for
all simplicial monomorphisms, by
\citet[Theorem I.7.10, pp.~32--35]{ABGoerssJardine} and the
boundary-by-boundary extension described in its proof. Apply this to
$L_C\hookrightarrow\Delta^{V(C)}$ to extend the displayed map to the
full simplex. The extension defines an object $D$ of
$\Delta_{\mathrm{inj}}/K$ on the same ordered vertex set. Every
previous copy of $A$ or $B$ remains an induced copy, because all of
its faces were already labelled and the extension preserves those
labels.

Now colour the arrows $A\to D$ with $c$ colours. Restrict the
colouring to the old copies $A\to C$. The Ramsey property of $C$
gives an old copy $B\to C$ whose $A$-subcopies are monochromatic.
This is a copy $B\to D$ with the same property. Hence $D$ is the
required Ramsey target. The category has finite hom-sets and joint
embeddings: the same extension argument applies to a disjoint union
of two labelled simplices. The integral affine-equalizer proof of
Ramsey exactness applies, and then works over every unital ring.
\end{proof}

\subsection{A counterexample to the all-prime implication}

\begin{lemma}[Proved: finite marginal lifting in a perfect group algebra]
\label{C4:perfect-marginals}
Let $p$ be a prime, let $V$ be an abelian group on which multiplication
by $p$ is bijective, and put $k=\mathbb F_p$. Suppose
$\rho_i:V\twoheadrightarrow V_i$ is a finite family of surjective
homomorphisms, with each $V_i$ uniquely $p$-divisible. Let
$J_i=\ker(k[V]\to k[V_i])$. Then the finite lattice of ideals generated
from the $J_i$ by sums and intersections is distributive. Consequently,
a family $z_i\in k[V_i]$ lifts to one $z\in k[V]$ if its images agree
in $k[V]/(J_i+J_j)$ for every pair $i,j$. The lift is a finite linear
combination of elements of $V$.

The same statement holds for affine surjections of torsors under these
groups: choose any origin in the source torsor, and the induced origins
in its quotients.
\end{lemma}
\begin{proof}
The Frobenius map on $k[V]$ sends $\sum a_v[v]$ to
$\sum a_v[pv]$. It is bijective, with inverse obtained by replacing
$v$ by its unique $p$th root. The same is true in every $k[V_i]$.
Thus each kernel $J_i$ is stable under Frobenius and its inverse.
Sums and intersections of ideals with this property retain it:
for sums use $(x+y)^{1/p}=x^{1/p}+y^{1/p}$, and for intersections
apply the inverse in each ideal. Every such ideal is radical. Indeed,
if $x^n$ belongs to it, then so does $x^{p^m}$ for $p^m\geq n$;
applying the inverse Frobenius $m$ times puts $x$ in the ideal.

The radical ideals of a commutative ring form a distributive lattice
with intersection as meet and radical of the sum as join. Here all
sums under consideration are already radical, so the operations are
ordinary sum and intersection. Equivalently, one can verify
\[
 (I+J)\cap(I+J')=I+(J\cap J')
\]
by taking radicals and using the corresponding identity of closed
sets in the prime spectrum. This proves the claimed distributivity.

For the lifting assertion induct on the number of ideals. Suppose
$x$ lifts the first $m-1$ residues and $y$ lifts the last. Pairwise
compatibility gives $x-y\in J_i+J_m$ for every $i<m$, whence
\[
 x-y\in\bigcap_{i<m}(J_i+J_m)
      =\left(\bigcap_{i<m}J_i\right)+J_m.
\]
Subtracting the first summand from $x$ makes it agree with $y$ modulo
$J_m$ without changing any earlier residue. All elements of a group
algebra have finite support. Changing an affine origin only translates
a basis, so the torsor formulation follows.
\end{proof}

\begin{lemma}[Proved: coherent modular weights on rational primitives]
\label{C4:primitive-weights}
Fix $r\geq2$ and a prime $p$. For every finite ordered rational
$r$-cocycle object $S$ there is a finitely supported element
\[
 \mu_S\in\mathbb F_p[\operatorname{Prim}(S)],\qquad
 \operatorname{Prim}(S)=
 \{b\in C^{r-1}(\Delta^S;\mathbb Q):db=c_S\},
\]
of augmentation one. These elements may be chosen simultaneously
so that every induced embedding $T\to S$ pushes $\mu_S$ to $\mu_T$
by restriction of primitives.
\end{lemma}
\begin{proof}
The primitive set is a nonempty affine torsor under the rational
vector space $V_S=Z^{r-1}(\Delta^S;\mathbb Q)$. Include the empty
structure for this construction only, with its unique primitive.
Restriction to a subset is surjective: a difference of two primitives
on that subset is a positive-degree cocycle, hence a coboundary;
extend its primitive cochain to $S$ and correct a chosen primitive
on $S$.

We need a precise two-subset assertion. If $T,U\subseteq S$, any
primitives on $T$ and $U$ agreeing on $T\cap U$ extend to a primitive
on $S$. Subtract the restriction of a fixed primitive on $S$. The
result is an $(r-1)$-cocycle on
$\Delta^T\cup\Delta^U$. That complex is contractible when its
intersection is nonempty and is a disjoint union of at most two
simplices otherwise. Its cohomology in degree $r-1>0$ vanishes.
Write the difference as a coboundary, extend the preceding cochain
to the full simplex, and add its coboundary to the fixed primitive.
This also proves, on the translation spaces,
\[
 \ker(V_S\to V_T)+\ker(V_S\to V_U)
   =\ker(V_S\to V_{T\cap U}).                       \tag{C4.1}
\]
For example, to prove the nontrivial inclusion, take $v$ vanishing
on the intersection and extend the compatible pair $(v|_T,0|_U)$
to $w\in V_S$; then $v=(v-w)+w$ is the required decomposition.
The kernels in the group algebra therefore satisfy
$J_T+J_U=J_{T\cap U}$: the kernel of a group-algebra quotient is
generated by $[h]-1$ for elements $h$ of its group kernel.

Choose the weights by induction on the size of $S$, one choice for
each ordered isomorphism type. At size zero choose the unique basis
element. Suppose the choices have been made naturally on all smaller
structures. For a structure $S$, prescribe the previously chosen
weight on each proper subset. These marginals agree after restriction
to every intersection, by the inductive hypothesis. Equation (C4.1)
and Lemma~\ref{C4:perfect-marginals} give an element of
$\mathbb F_p[\operatorname{Prim}(S)]$ with all the prescribed
marginals. Include the empty marginal, so its augmentation is one.
A finite linearly ordered structure has no nonidentity increasing
self-isomorphism. Thus transporting this choice from the representative
to every ordered copy is unambiguous, and supplies all embedding
compatibilities. This completes the induction. The argument gives
finite support at every individual object, although no uniform bound
on support sizes is asserted.
\end{proof}

\begin{theorem}[Refuted: all-prime exactness forces rational exactness]
\label{C4:counterexample}
For each $r\geq2$, the category $A=\mathcal A_r(\mathbb Q)$ of finite
nonempty ordered rational $r$-cocycles is exact over $\mathbb F_p$
for every prime $p$, but is not exact over $\mathbb Q$ or $\mathbb Z$.
It has finite hom-sets and common successors. It lies in setting (S)
and satisfies (Rig), (SA), (FA), (St), (Op), (WEI), and (Fin).
It is not total, since $BA\simeq K(\mathbb Q,r)$.

Consequently the main all-prime implication in C.4 is false, both
with and without the finite-hom hypothesis, and remains false in
setting (S) with (Rig). Property (P) does not imply totality, even
for these nonempty connected finite-hom rigid strong Fra\"iss\'e
categories. The version whose index category is assumed total is
not decided by this example.
\end{theorem}
\begin{proof}
Fix $p$. The expansion category $B$ consists of finite ordered
$(r-1)$-cochains with rational values. Its reduct sends $b$ to $db$.
For each finite collection of objects and arrows, only finitely many
cochain values occur. Encode these values by finitely many relation
symbols, with value zero represented by absence of a relation.
The corresponding class of coloured $r$-uniform hypergraphs has free
amalgamation, so its arbitrary linear-order expansion is Ramsey by
\cite[Theorem~1.3]{D-EHN2021}. Hence every finite family of parallel
arrows in $B$ has an integral affine equalizer, by the rechecked
Ramsey-to-affine argument. Enlarging the finite set of colours to
include zero preserves all required objects and embeddings.

We give the finite-family passage explicitly, so no inference from
constant-coefficient or action-component homology is needed. Let
$a_0,\ldots,a_m:P\to Q$ be any finite parallel family in $A$.
Write the finite support of Lemma~\ref{C4:primitive-weights} as
\[
 \mu_Q=\sum_{b\in\mathcal S}w_b[b],\qquad
 \sum_b w_b=1\quad\text{in }\mathbb F_p.
\]
Jointly embed the finitely many primitive expansions $Q_b=(Q,b)$
into one cochain object $X^*$, by maps $j_b:Q_b\to X^*$.
For every primitive expansion $P_d$ occurring as some
$a_i^*b$, consider all the maps
\[
 j_b a_i:P_d\longrightarrow X^*
       \quad\text{with }a_i^*b=d.
\]
There are finitely many such types $d$ and finitely many such maps.
Use integral affine equalizers in $B$ successively for these families,
postcomposing later families with the previously obtained affine
arrows. Equivalently, at a later stage equalize the finitely many
ordinary arrows in the supports of the corresponding affine arrows.
This yields one affine arrow $s:X^*\to Y^*$ over $\mathbb F_p$
for which all the displayed maps of each fixed type $d$ have the
same image $v_d$. Postcomposition preserves every equality already
achieved.

Let $u:B\to A$ forget the primitive and set
\[
 t=\sum_{b\in\mathcal S}w_b\,u(s)u(j_b)
      \in\mathbb F_p[A(Q,uY^*)].
\]
Its augmentation is one. For every $i$ its composite with $a_i$ is
\[
 ta_i=\sum_d\left(\sum_{b:a_i^*b=d}w_b\right)u(v_d)
      =\sum_d\mu_P(d)\,u(v_d),
\]
independent of $i$, by coherence of the weights. Thus every finite
parallel family has an affine equalizer. The affine criterion proves
exactness over $\mathbb F_p$ for arbitrary module diagrams.

For rational nonexactness one can use a finite named family, not just
the nonzero homology of the base. On $Q=\{0,\ldots,2r+1\}$ choose an
$(r-1)$-cochain taking value $1$ on $(1,\ldots,r)$ and zero on every
other increasing $r$-tuple, and give $Q$ its coboundary cocycle.
The first $(r+1)$-point block has cocycle value one; the disjoint
second block has value zero. Let $P$ have $r$ points and zero
$r$-cocycle, and take the $2(r+1)$ embeddings $P\to Q$ onto the
codimension-one faces of these two blocks. At any target $D$, choose
a rational primitive $b_D$. If an affine combination $t$ equalized
all these face embeddings, evaluating $b_D$ and alternating over
the faces of the two blocks would give equal answers. The coboundary
identity instead gives $\operatorname{aug}(t)=1$ on the first block
and zero on the second. This contradiction excludes a rational
affine equalizer at every target. Clearing denominators also excludes
integral exactness.

The strong-class properties and the identification
$BA\simeq K(\mathbb Q,r)$ are the audited constructions of R13:
all subsets are induced strong substructures, cochains extend across
two simplices, nonempty-base amalgamation is functorial, and the
three-piece extension uses $H^r(S^1;\mathbb Q)=0$ for $r\geq2$.
The minimal finite-support argument gives (Op), and rigidity comes
from the order. For (Fin), if $G_C\leq G_B$ has finite index, every
point of $C$ has finite $G_B$-orbit. Free duplication over $B$ forces
$C\subseteq B$, hence $G_B\leq G_C$ and equality. (WEI) follows
from (Op). Finite hom-sets are subsets of sets of increasing
injections between finite ordered sets. Joint embeddings give
common successors. Finally $H_r(BA;\mathbb Z)\cong\mathbb Q\ne0$,
so $BA$ is not contractible and $A$ is not total. The established
profinite/all-prime equivalence gives (P).
\end{proof}

\begin{corollary}[Proved: the complete field profile of cocycle classes]
\label{C4:field-profile}
Let $\Lambda$ be a countable abelian group and let $r\geq2$. Then
\[
 \mathcal A_r(\Lambda)\text{ is exact over }\mathbb F_p
 \quad\Longleftrightarrow\quad
 p:\Lambda\longrightarrow\Lambda\text{ is bijective},
\]
and
\[
 \mathcal A_r(\Lambda)\text{ is exact over }\mathbb Q
 \quad\Longleftrightarrow\quad \Lambda\text{ is torsion}.
\]
Thus it is exact over every prime field $\mathbb F_p$ precisely when
$\Lambda$ is a rational vector space, and it is integrally exact
precisely when $\Lambda=0$. For a nonzero rational vector space these
categories give the counterexamples in Theorem~\ref{C4:counterexample}.
\end{corollary}
\begin{proof}
If multiplication by $p$ is bijective on $\Lambda$, the primitive
translation group
$Z^{r-1}(\Delta^S;\Lambda)$ is uniquely $p$-divisible. This follows
either from the cone contraction, which identifies it with a finite
power of $\Lambda$, or by applying the inverse of multiplication by
$p$ to cochains coordinatewise. The two-subset extension and the
perfect group-algebra argument of
Lemmas~\ref{C4:perfect-marginals}--\ref{C4:primitive-weights} therefore
apply verbatim. The finite-family proof then gives exactness over
$\mathbb F_p$ for every module diagram.

For necessity, exactness annihilates all positive derived colimits,
in particular homology with the constant $\mathbb F_p$ diagram.
The base has $H_r(BA;\mathbb Z)=\Lambda$, and no reduced integral
homology below degree $r$. The universal coefficient theorem gives
\[
 H_r(BA;\mathbb F_p)=\Lambda/p\Lambda,
 \qquad
 H_{r+1}(BA;\mathbb F_p)\twoheadrightarrow\Lambda[p].
\]
Both groups on the right must vanish. Thus multiplication by $p$ is
surjective and injective. This use of constant coefficients proves
necessity only; sufficiency was established above for arbitrary
modules.

If $\Lambda$ is torsion, the finitely many values occurring in a
fixed finite parallel family generate a finite subgroup
$\Lambda_0$. The finite-group primitive-averaging proof in
Theorem~\ref{A:finite-cocycles} supplies a rational affine equalizer
inside $\mathcal A_r(\Lambda_0)$, hence in $\mathcal A_r(\Lambda)$.
The affine criterion proves rational exactness. Conversely, if
$\Lambda$ is not torsion, choose a nonzero linear functional on
$\mathbb Q\otimes\Lambda$ and normalize it at a nonzero image of
some $\lambda\in\Lambda$. It induces a homomorphism
$\chi:\Lambda\to\mathbb Q$ with $\chi(\lambda)=1$.
Apply $\chi$ to a target primitive in the finite two-block face
obstruction of Theorem~\ref{C4:counterexample}. The same alternating
sum yields $1=0$, so rational exactness fails.

Multiplication by every prime is bijective exactly for rational
vector spaces. Integral exactness would force the constant homology
$H_r(BA;\mathbb Z)=\Lambda$ to vanish; for $\Lambda=0$ the category
is $\Delta_{\mathrm{inj}}$, which is integrally exact. This proves
the remaining assertions.
\end{proof}

\begin{example}[A finite modular primitive distribution]
\label{C4:triangle-distribution}
For a triangle of cocycle value one, write a primitive as
$(x,y,z)=(b_{01},b_{12},b_{02})$, so $x+y-z=1$.
For any prime $p$, put weight $+1$ on
\[
 (1/p,j/p,(j+1-p)/p),\qquad 0\leq j<p,
\]
and weight $-1$ on
\[
 (0,j/p,(j-p)/p),\qquad 1\leq j<p.
\]
This has $2p-1$ distinct atoms and integral augmentation one. Its
three marginals are respectively
$p[1/p]-(p-1)[0]$, $[0]$, and $[0]$.
Modulo $p$ all are $[0]$. This exhibits explicitly how rational
denominators create the modular marginal cancellation, while the
rational cocycle still has value one. The exact-arithmetic program
\texttt{C4\_perfect\_marginal\_certificate.py} checks the formulas
for eight primes; the displayed calculation proves them uniformly.
\end{example}

\begin{proposition}[Proved: a bounded-denominator criterion]
\label{C4:noetherian}
Let $A$ be nonempty with common successors, and fix a finite parallel
family $a_0,\ldots,a_k:p\to q$. Write
\[
 D:\mathbb Z[A(q,-)]\longrightarrow
       \mathbb Z[A(p,-)]^k,
 \qquad D(t)=(ta_i-ta_0)_{1\leq i\leq k},
\]
and let $S$ be the pointwise saturation of $\operatorname{im}D$ in
the target. If $S/\operatorname{im}D$ is annihilated by a positive
integer $N$, an affine equalizer over $\mathbb F_p$ for any prime
$p\nmid N$ implies a rational affine equalizer for this family.
Consequently all-prime exactness implies rational exactness whenever
these torsion quotients have bounded exponent for every finite family.

The bounded-exponent condition holds if each such saturation functor
is generated by finitely many elements at finitely many objects;
in particular it holds if finite sums of integral representable
functors are Noetherian. In setting (S), an alternative sufficient
condition is that every finite sum of the smooth permutation modules
$\mathbb Z[G/G_P]$ is Noetherian as a $\mathbb Z[G]$-module.
No Noetherian hypothesis is asserted for arbitrary oligomorphic
groups or for the independently ordered vector-space example.
\end{proposition}
\begin{proof}
Choose a modular affine equalizer at a target $r$ and lift it to
$t\in\mathbb Z[A(q,r)]$ with $\operatorname{aug}(t)\equiv1\pmod p$.
Then $D(t)=p y$ for some integral vector $y$. The class of $y$ in
$S(r)/\operatorname{im}D(r)$ is killed by both $p$ and $N$, hence is
zero. Write $y=D(z)$ with $z$ integral. Now
$w=t-pz$ lies in $\ker D(r)$ and its augmentation is congruent to
one modulo $p$, so is a nonzero integer. Dividing $w$ by that
augmentation gives the required rational affine equalizer.

Suppose $S$ is generated by $s_j\in S(r_j)$ for $1\leq j\leq m$.
For each $j$ choose $n_j>0$ with $n_js_j\in\operatorname{im}D(r_j)$.
The product $N=\prod_jn_j$ then annihilates the saturation quotient
at every object, by naturality and generation. A subfunctor of a
Noetherian finite sum of representables is finitely generated, which
proves the categorical sufficient condition.

For the group-theoretic sufficient condition, take the filtered union
over finite strong subsets of $M$. The resulting map is
\[
 D_\infty:\mathbb Z[\operatorname{Emb}(q,M)]\longrightarrow
             \mathbb Z[\operatorname{Emb}(p,M)]^k.
\]
Its saturated image is a submodule of the indicated Noetherian
permutation module, hence has finitely many $\mathbb Z[G]$-generators.
An integer multiple of each generator belongs to
$\operatorname{im}D_\infty$, so a common positive integer annihilates
the saturation quotient. The preceding modular lifting argument
works for $D_\infty$. All finitely many embeddings and coefficients
which occur in the resulting rational vector lie in one finite
strong subset of $M$, so it is a rational affine equalizer in $A$.
The all-family affine criterion finishes the proof.
\end{proof}

\paragraph{Remaining total-category gate in C.4.}
The counterexamples above have noncontractible base and hence do not
decide whether all-prime exactness implies rational exactness when
$A$ is additionally assumed total. The arithmetic conclusions of
Proposition~\ref{C:rational-gate} still apply to that restricted
question. Theorem~\ref{C4:contractible-kan-ramsey} excludes a tempting
candidate family: replacing $K(\mathbb Q,r)$ by a contractible Kan
complex in $\Delta_{\mathrm{inj}}/K$ always gives an integrally exact
Ramsey category. The stronger
Theorem~\ref{C:kan-filtered-totalization} excludes every filtered
diagram of connected Kan simplex categories with weakly contractible
colimit base, allowing arbitrary transition functors and arbitrary
coefficient groups. Such a colimit is always integrally exact.
Theorem~\ref{C4:acyclic-cocycle-colimits} also rules out rationally
acyclic filtered colimits of rational-vector coefficient cocycle
categories, without requiring contractibility as a hypothesis.
More generally, Theorem~\ref{C4:exact-base-kan-assembly} excludes
every total Grothendieck assembly of connected Kan simplex
categories over an integrally exact base, even when the base is
not filtered and the transition functors are arbitrary.
The coefficient-pushforward special case in
Theorem~\ref{C4:grothendieck-exclusion} is total precisely when
its coefficient diagram is locally zero.
No total counterexample or proof of the restricted implication is
asserted here.

\begin{corollary}[Proved: a compact torsor witnessing the all-prime failure]
\label{C4:compact-obstruction}
For $A=\mathcal A_2(\mathbb Q)$ there is a diagram $C:A\to\Ab$
that is finitely generated at each object and satisfies
\[
 T_1(C)\cong\mathbb Z,
 \qquad
 \lim_{A^{\mathrm{op}}}^{1} C^\vee\cong\mathbb R/\mathbb Z.
\]
Every profinite abelian diagram on $A^{\mathrm{op}}$ has vanishing
first derived limit, while the displayed compact diagram does not.
A nonzero element of $\mathbb R/\mathbb Z$ gives an objectwise
nonempty compact torsor diagram with empty inverse limit.
\end{corollary}
\begin{proof}
Use the six face embeddings $a_0,\ldots,a_5:P\to Q$ in
Theorem~\ref{C4:counterexample}, and set
\[
 D:\mathbb Z[A(Q,-)]\longrightarrow\mathbb Z[A(P,-)]^5,
 \qquad D(t)=(ta_i-ta_0)_{1\leq i\leq5},
 \qquad C=\operatorname{coker}D.
\]
The rational primitive functional in that proof shows
$\operatorname{aug}\ker D(X)=0$ at every object $X$.
Writing $J=\operatorname{im}D$, right exactness gives
$\operatorname{colim}J\cong\mathbb Z$. Its map to the colimit
$\mathbb Z^5$ of the target representable sum is zero. That sum
is projective, so the long exact sequence of
$0\to J\to\mathbb Z[A(P,-)]^5\to C\to0$ gives
$T_1(C)\cong\mathbb Z$. This argument does not require $D$ itself
to be injective. Finite hom-sets make every $C(X)$ finitely generated.
The duality theorem in Part C gives the asserted first derived limit.

For a direct torsor description, fix $0\ne\theta\in\mathbb R/\mathbb Z$.
Inside
$((\mathbb R/\mathbb Z)^{A(P,X)})^5$ take the affine equations
\[
 \sum_{i=1}^5\bigl(z_i(ga_i)-z_i(ga_0)\bigr)=\theta
       \qquad(g\in A(Q,X)).
\]
Their homogeneous solution group is exactly $C(X)^\vee$.
They are solvable: the homomorphism
$t\mapsto\theta\operatorname{aug}(t)$ on
$\mathbb Z[A(Q,X)]$ vanishes on $\ker D(X)$, hence descends to
$J(X)$ and extends to the ambient free abelian group because
$\mathbb R/\mathbb Z$ is injective. Restrictions along embeddings
preserve the equations. A compatible solution would make each
coordinate function $z_i$ constant on embeddings of $P$, by
restriction along those embeddings from its value at $\operatorname{id}_P$.
The left side of every equation would then be zero, contradicting
$\theta\ne0$. Thus the torsor limit is empty. Finally all-prime
exactness and the profinite test theorem imply vanishing on every
profinite abelian diagram, as claimed.
\end{proof}


\subsection{Explicit compact obstructions over \texorpdfstring{$\mathcal C_\delta$}{C-delta}}

\begin{proposition}[C.8; Proved: the torsor and its nonzero cocycle]\label{C:explicit-torsor}
Let $A=\Cdelta$, $P_X=A(V,X)$, $E_X=A(W,X)$, and let $k=\F_p$. Define $D_X:k[E_X]\to k[P_X]^2$ by
\[
 D_X[g]=([gv_2]-[gv_1],\ [gv_3]-[gv_1]),\qquad Q=\coker D.
\]
The following contravariant diagram consists of finite-dimensional $k$-vector spaces:
\[
 K(X)=\left\{(s,t)\in k^{P_X}\times k^{P_X}:
 s(gv_2)-s(gv_1)+t(gv_3)-t(gv_1)=0\ (g\in E_X)\right\}.
\]
Its transitions are restriction along postcomposition of embeddings. It is $\Hom_k(Q,k)$ and satisfies $\lim^1K\cong k$. Replacing the right-hand side zero in its defining equations by one gives a nonempty torsor $T(X)$ under $K(X)$ at each object, and $\lim T=\varnothing$.

Over $\Z$, use the same $D$, $Q$ and equations with values in $\mathbb T$. Then $K=\Hom_{\Z}(Q,\mathbb T)$ and $\lim^1K\cong\mathbb T$. For each nonzero $\theta\in\mathbb T$, the equations with right-hand side $\theta$ give a torsor without a compatible point.
\end{proposition}
\begin{proof}
We recheck the input from T~A.24 rather than assuming it. In the support of a nonzero element of $k[E_X]$, different increasing embeddings have different hyperedge images. If there are $r$ such images on $v$ vertices, positivity of $2v-3|E_X(U)|$ gives $2v>3r$. Therefore some vertex belongs to exactly one image. Its coefficient in at least one of the two displayed difference coordinates is the original nonzero coefficient, up to sign. Thus $D_X$ is injective, including in characteristic two. The same argument works over $\Z$. The two-term projective resolution
\[
 0\longrightarrow P_W\xrightarrow D P_V^2\longrightarrow Q\longrightarrow0
\]
has colimit differential zero. Hence $T_1(Q)=k$ or $\Z$, respectively. Theorem~\ref{C:duality} gives the stated derived limits and also shows that the duals of $P_W$ and $P_V^2$ have no positive derived limits.

Dualizing the injection $D_X$ gives a surjection onto $k^{E_X}$ or $\mathbb T^{E_X}$, so the indicated fibers are nonempty. Restriction carries fibers to fibers because it preserves the equations. If a compatible family existed, its two functions on $P_X$ would be constant: for each $v:V\to X$, compatibility identifies their value at $v$ with their value at $\id_V$. The equation for $\id_W$ would then give $0=1$ or $0=\theta$, a contradiction.

For the actual cocycle, choose any lift $z_X=(s_X,t_X)$ of the constant function $1$ (or $\theta$). In the finite-field case a completely specified choice is the first solution in lexicographic order using the skeleton's ordered hom-sets. For $f:X\to Y$ put
\[
 c_f=z_Y|_f-z_X\in K(X).
\]
Then $c_{gf}=c_f+K(f)c_g$. The connecting map for
$0\to K\to P_V^{\vee\,2}\to P_W^\vee\to0$
identifies its class with $1\in k$ or $\theta\in\mathbb T$, since the map of ordinary limits $k^2\to k$ (respectively $\mathbb T^2\to\mathbb T$) is zero. This both computes the class and proves it is nonzero.
\end{proof}

\begin{theorem}[C.9; Refuted: objectwise dimensions determine the first derived limit]\label{C:dimensions}
For every prime $p$ there are two diagrams $F,F':\Cdelta^{\op}\to\F_p\text{-}\mathrm{Mod}$, finite dimensional at every object, with
\[
 \dim F(X)=\dim F'(X)\quad\text{for all }X,
 \qquad \dim\lim^1F=1,\quad \dim\lim^1F'=0.
\]
More generally, for any small category $A$, any unital ring $R$, and any objectwise injective map
\[
 D:\bigoplus_{j=1}^{m}R[A(p_j,-)]\longrightarrow
       \bigoplus_{i=1}^{n}R[A(q_i,-)],
\]
let $B:R^m\to R^n$ be its induced augmentation map. Then
\[
 T_1(\coker D)=\ker B,\qquad
 T_i(\coker D)=0\ (i\geq2),\qquad
 \colim_A\coker D=\coker B.
\]
Here $B$ is a morphism of left $R$-modules; for a noncommutative $R$ its entries act on a coordinate on the right.
\end{theorem}
\begin{proof}
Take $F=\Hom_k(Q,k)$ from Proposition~\ref{C:explicit-torsor}. Choose $(\alpha,\beta,\gamma)=(1,1,1)$ if $p\ne3$ and $(1,1,2)$ if $p=3$. All three entries and their sum are nonzero. Define
\[
 D'':k[A(W,-)]\to k[A(V,-)],\qquad
 [g]\mapsto\alpha[gv_1]+\beta[gv_2]+\gamma[gv_3].
\]
The private-vertex proof makes it injective. Its augmentation is the nonzero scalar $\alpha+\beta+\gamma$, so its cokernel has zero first derived colimit. Put
\[
 Q'=k[A(V,-)]\oplus\coker D'',\qquad F'=\Hom_k(Q',k).
\]
Both $Q(X)$ and $Q'(X)$ have dimension $2|P_X|-|E_X|$. The above projective resolutions and duality give the claimed derived-limit dimensions. The cocycle of the nonzero class is exactly $c_f$ in Proposition~\ref{C:explicit-torsor}; it is attached to $F$, while no nonzero class exists for $F'$.

For the general formula, the assumed injection followed by its cokernel is a projective resolution of length one. The colimit of each representable free module is $R$, with its basis embeddings all identified: the element category $p/A$ has initial object $\id_p$. If
\[
 D_{p_j}[\id_{p_j}]=\sum_{i,h:q_i\to p_j}b_{ijh}[h],
\]
then $B$ sends $(r_j)_j$ to $(\sum_{j,h}r_jb_{ijh})_i$. Applying colimit to the resolution proves every formula. In particular sizes of the objectwise matrices do not determine the kernel of their augmentation map.
\end{proof}

\subsection{The four-object exit category}

\begin{theorem}[C.10; Proved: the exit index and its constant coefficients]\label{C:exit}
Let $G=F(a,b)$, let $c=(ab)^{-1}$, and put $H_0=\langle a\rangle$, $H_1=\langle b\rangle$, $H_\infty=\langle c\rangle$. Define a category $X$ on the objects $0,1,\infty,U$ by
\[
 X(U,U)=G,\quad X(x,U)=G/H_x\ (x=0,1,\infty),\quad
 X(x,x)=\{\id_x\},
\]
with all other hom-sets empty. Here $G/H_x$ means left cosets $gH_x$, and composition is $h\circ(gH_x)=hgH_x$ and $h\circ g=hg$. This is a skeleton of the homotopy exit-path category of $\mathbb P^1$ stratified by its three punctures and their complement. Its classifying space is homotopy equivalent to $S^2$. Consequently
\[
 T_1^X(\Delta\Z)=0,\quad T_2^X(\Delta\Z)=\Z,
 \qquad \lim^1_{X^{\op}}\Delta\mathbb T=0,\quad
 \lim^2_{X^{\op}}\Delta\mathbb T=\mathbb T.
\]
\end{theorem}
\begin{proof}
Choose a base point of the open stratum and disjoint access paths to small circles around the punctures. The open stratum is a punctured sphere with fundamental group $F(a,b)$ and peripheral elements $a,b,c$, with $abc=1$. An exit path from a puncture to the base point is determined by a path in the open stratum modulo the peripheral loop at its starting end. These are exactly the cosets displayed above; concatenation at the terminal base point is left multiplication. A different system of access paths conjugates the peripheral subgroups and gives an isomorphic category. The three point strata have trivial endomorphism groups and there are no exit paths to a lower stratum.

This category also models the full exit-path homotopy type. The open stratum is aspherical, and each peripheral circle is aspherical and injects on fundamental groups. The homotopy fibers describing its exit mapping spaces therefore have contractible components. Thus replacing those mapping spaces by their sets of components loses no homotopy information.

More concretely, adjoining a source object with arrow set $G/H$ to the one-object group category $G$ produces the homotopy cofiber of $BH\to BG$. Indeed the part of its nerve consisting of simplices that start at that source is the cone on the nerve of the action groupoid $G\ltimes(G/H)$; the latter is equivalent to $BH$. With three source objects the classifying space is the homotopy pushout attaching three such cones. Since $BG$ is a wedge of two circles and the three attaching loops are $a,b,(ab)^{-1}$, this is the presentation two-complex with those three relators. The first two two-cells collapse the two circles, and the remaining cell is attached by a nullhomotopic loop. The resulting space is $S^2$. The homology assertions follow, and Theorem~\ref{C:duality} gives the derived-limit assertions. We give the actual bar cycle and cobar cocycle next, rather than leaving this last computation solely to duality.
\end{proof}

\begin{proposition}[C.11; Proved: an explicit degree-two cobar class]\label{C:exit-cocycle}
Use the category and composition convention of Theorem~\ref{C:exit}. For each $x$ choose a coset section $s_x:G/H_x\to G$ with $s_x(H_x)=1$. Define the integer $n_x(g,z)$ by
\[
 s_x(gz)^{-1}g s_x(z)=\gamma_x^{\,n_x(g,z)},
 \qquad (\gamma_0,\gamma_1,\gamma_\infty)=(a,b,c).
\]
For each $\theta\in\mathbb T$ there is a normalized degree-two cocycle $\omega_\theta$ with value
$\theta n_\infty(g,z)$ on $\infty\xrightarrow zU\xrightarrow gU$, and value zero on all other nondegenerate two-simplices. The map
\[
 \mathbb T\longrightarrow\lim^2_{X^{\op}}\Delta\mathbb T,
 \qquad \theta\longmapsto[\omega_\theta]
\]
is an isomorphism. There is no single group-theoretic generator of $\mathbb T$; the requested ``generating cocycle'' is therefore properly a family indexed by $\theta$, or the integral generator whose scalar multiples are this family.
\end{proposition}
\begin{proof}
The section identity gives
\[
 n_x(hg,z)=n_x(h,gz)+n_x(g,z).
\]
For a three-simplex starting at $x$, this is exactly the cobar cocycle equation; on three-simplices inside $G$ all values are zero. Thus $\omega_\theta$ is a cocycle.

For an explicit integral cycle write $u_x=H_x:x\to U$ and use $[f,g]$ for the simplex with first arrow $f$, then arrow $g$. In the normalized bar complex put
\[
 z=\sum_{x=0,1,\infty}[u_x,\gamma_x]
       -[b,a]-[c,ab].
\]
We have $\partial[u_x,\gamma_x]=[\gamma_x]$, because $\gamma_xu_x=u_x$. Also
\[
 \partial([b,a]+[c,ab])=[a]+[b]+[c],
\]
since $(ab)c=1$ and identity edges are degenerate. Hence $\partial z=0$. In the cone-attachment model, this is the fundamental cycle obtained from the three peripheral disks and the fundamental relation $abc=1$. It generates $H_2(X;\Z)=\Z$. Direct evaluation gives $\omega_\theta(z)=\theta$.

For a direct cohomological description, the cone-attachment cobar complex has its degree-two quotient
\[
 \mathbb T^3/\{(r,s,-r-s):r,s\in\mathbb T\}\cong\mathbb T,
\]
with isomorphism the sum of coordinates. Its degree-one kernel is zero. The vector $(0,0,\theta)$ is represented in the full cobar complex by $\omega_\theta$, and it maps to $\theta$ under this quotient. This verifies both derived limits and their agreement with the bar duality. The code certificate checks the displayed boundary identity in freely reduced words and the integral cellular boundary matrix.
\end{proof}

\begin{proposition}[C.12; Proved: nonconstant coefficients detect degree one]\label{C:exit-nonconstant}
The category $X$ is not quasi-filtered and is not exact over any nonzero ring. It satisfies none of (CH), (CC), (CT). For its constant compact coefficients the first nonzero positive derived limit occurs in degree two; for an explicit nonconstant compact abelian diagram it already occurs in degree one.
\end{proposition}
\begin{proof}
Let $P=R[X(U,-)]$. It is zero at the three point objects and is the regular left $R[G]$-module at $U$. Define the natural map
\[
 D:P\to P,\qquad [g]\mapsto[ga]-[g].
\]
It is objectwise injective over every nonzero ring: the orbits of right multiplication by the infinite-order element $a$ are copies of $\Z$, on which a finitely supported invariant function is zero. This proof needs no domain hypothesis and no zero-divisor theorem for group rings. The induced map on colimits is zero, so for $M=\coker D$ the representable resolution gives $T_1(M)=R$. The identical finite-support argument rules out an affine equalizer of $\id_U,a$, proving failure of IV.

For $R=\Z$, the compact dual $K=M^\vee$ is the one-point group at each puncture and $\mathbb T^{G/H_0}$ at $U$, with the induced permutation action. It has $\lim^1K=\mathbb T$. For $0\ne\theta\in\mathbb T$ its nontrivial torsor is a point at the punctures and
\[
 T(U)=\{f\in\mathbb T^G:f(ga)-f(g)=\theta\text{ for all }g\in G\}.
\]
This set is nonempty by choosing values on the right $\langle a\rangle$-cosets and extending along their integer coordinates. The action maps come from the contravariant dual of left translation, and the maps to the point objects are unique. A compatible point would be a translation-invariant $f$, which would contradict its defining equation. Thus (CT) and (CH) fail.

For (CC), take at $U$ the compact convex space of probability measures on the boundary of the free group $F(a,b)$, and at each puncture a singleton. The boundary is the compact space of infinite reduced words in $a^{\pm1},b^{\pm1}$; use inverse pushforward for the contravariant group action. For each letter $s$ let $C_s$ be the clopen cylinder of words starting with $s$. Left multiplication by $a$ maps the complement of $C_{a^{-1}}$ homeomorphically onto $C_a$, and similarly for $b$. An invariant probability measure would satisfy
\[
 \mu(C_a)+\mu(C_{a^{-1}})=1,
 \qquad \mu(C_b)+\mu(C_{b^{-1}})=1.
\]
The four disjoint cylinders partition the boundary, so their total mass would be both one and two. Hence the diagram has empty limit. Its spaces are compact convex subsets of the continuous dual of the real Banach space of continuous boundary functions, with the weak-star topology, and its transition maps are continuous affine maps. This establishes (CC)'s failure without merely quoting nonamenability.
\end{proof}

% Further coefficient reductions, appended without renumbering earlier statements.
% Appended to Part C to preserve previously assigned statement numbers.
\begin{proposition}[Proved: a finite rational primitive criterion]
\label{C4:rational-primitive-criterion}
Let $s\geq2$ and let $a_0,\ldots,a_k:P\to Q$ be a finite family
in $\mathcal A_s(\mathbb Q)$. The following are equivalent:
\begin{enumerate}
\item the family has an affine equalizer over $\mathbb Q$;
\item there are primitives $b_P\in C^{s-1}(\Delta^P;\mathbb Q)$
and $b_Q\in C^{s-1}(\Delta^Q;\mathbb Q)$ of the specified cocycles
such that $a_i^*b_Q=b_P$ for every $i$;
\item the family has an affine equalizer over $\mathbb Z$.
\end{enumerate}
Thus rational equalizability for a fixed family is a finite system
of rational linear equations on its two given objects. This
assertion concerns the coefficient group $\mathbb Q$ of the
cocycles; it is not asserted with a finite coefficient group in its
place.
\end{proposition}
\begin{proof}
Suppose $t=\sum_g\lambda_g g\in
\mathbb Q[\mathcal A_s(\mathbb Q)(Q,D)]$ is an affine equalizer.
Choose a primitive $b_D$ of the cocycle on the simplex $D$ and set
\[
 b_Q=\sum_g\lambda_g g^*b_D.
\]
Because $\sum_g\lambda_g=1$, its coboundary is the cocycle on
$Q$. The equality of the formal composites $ta_i$ makes all
$a_i^*b_Q$ equal; their common value is the required $b_P$.
Conversely such primitives lift the whole family to one pair of
objects in the primitive-cochain expansion. That expansion is the
Ramsey class of ordered arbitrary $(s-1)$-cochains used in
Theorem~\ref{C4:counterexample}. Its integral affine criterion gives
an integral equalizer of the lifted family, and forgetting the
primitives gives one in $\mathcal A_s(\mathbb Q)$. Finally an
integral equalizer is a rational one. The converse uses the fact
that the barycenters above still take values in the actual cocycle
coefficient group $\mathbb Q$.
\end{proof}

\begin{theorem}[Proved: increasing cocycle degrees erase fixed families]
\label{C4:degree-colimit-exclusion}
Let $r,s\geq2$ and let
$F:\mathcal A_r(\mathbb Q)\to\mathcal A_s(\mathbb Q)$ be any
functor. If $a_0,\ldots,a_k:P\to Q$ is a finite parallel family
and $s\geq |Q|$, the family $F(a_0),\ldots,F(a_k)$ has an integral
affine equalizer. In particular, for any sequence of functors
\[
 \mathcal A_{r_0}(\mathbb Q)\longrightarrow
 \mathcal A_{r_1}(\mathbb Q)\longrightarrow\cdots,
 \qquad r_n\longrightarrow\infty,
\]
its filtered colimit as a small category is integrally exact.
Consequently this construction cannot give a total all-prime-exact
category which fails rational exactness, regardless of the chosen
transition functors.
\end{theorem}
\begin{proof}
We work with the equivalent category of finite subsets of the
Fraisse limit $M$ of $\mathcal A_r(\mathbb Q)$ and write
$G=\operatorname{Aut}(M)$. The previously checked cocycle-class
properties give (Op) and
$\operatorname{Fix}_M(G_B)=B$ for every finite $B\subset M$.
All finite subsets are strong, and the points of $M$ form one
$G$-orbit.

Take the filtered colimit of the underlying ordered sets of $F(B)$
over inclusions of finite nonempty subsets $B\subset M$. It is a
nonempty ordered smooth $G$-set $S$, since every transition map is
an increasing injection. Its $G$-action is induced by
$F(g:B\to gB)$. The cocycles on the $F(B)$ agree under inclusions,
so give a $G$-invariant cocycle
\[
 c\in C^s(\Delta^S;\mathbb Q).
\]
Here a finite tuple of vertices of $S$ lies in some $F(B)$, so this
definition covers every simplex and its coboundary equation.

If $S$ has a $G$-fixed vertex, coning chains to that vertex gives
a $G$-invariant primitive of $c$ on the entire simplex. Its
restrictions to $F(P)$ and $F(Q)$ satisfy the preceding proposition,
and we are done. Suppose therefore that $S$ has no fixed vertex.
For each $v\in S$, (Op) writes its stabilizer uniquely as
$G_{B_v}$ for a finite nonempty subset $B_v\subset M$. Uniqueness
follows from the displayed fixed-point identity. Moreover
\[
 B_{gv}=gB_v,
 \qquad v\in F(B)\ \Longrightarrow\ B_v\subset B,
 \tag{DC.1}
\]
because $G_B\subset G_v=G_{B_v}$ in the second assertion.
Choose $x\in M$ and a vertex $v\in F(\{x\})$. By (DC.1),
$B_v=\{x\}$, so its orbit supplies a $G$-equivariant injection
$j:M\to S$ with $j(x)=v$.

On rational vertex spaces define the augmentation-preserving,
$G$-equivariant maps
\[
 p:\mathbb Q[S]\longrightarrow\mathbb Q[M],
 \quad p(v)=\frac1{|B_v|}\sum_{x\in B_v}x,
 \qquad h=jp:\mathbb Q[S]\longrightarrow\mathbb Q[S].
\]
Use oriented simplex chains, identifying
$C_n(\Delta^S;\mathbb Q)$ with
$\bigwedge^{n+1}\mathbb Q[S]$ and the boundary with contraction
by augmentation. On each increasing basis tuple put
\[
\begin{split}
 H_n(v_0\wedge\cdots\wedge v_n)
 =\sum_{i=0}^n(-1)^i\,&h(v_0)\wedge\cdots\wedge h(v_i)\wedge v_i
 \\[-2pt]
 &{}\wedge\cdots\wedge v_n .
\end{split}
\tag{DC.2}
\]
Extend this formula linearly on the increasing-tuple basis. All
terms have finite support. The ordinary prism telescoping calculation
gives
\[
 \partial H+H\partial=\operatorname{id}-\bigwedge h.
 \tag{DC.3}
\]
For example, in degree zero it is
$\partial(h(v)\wedge v)=v-h(v)$; in higher degrees the internal
faces cancel in consecutive summands. The same calculation is valid
for rational barycenters because every $h(v)$ has augmentation one.
The use of increasing tuples is essential to the definition:
$G$ preserves their order, and hence (DC.2) is $G$-equivariant.
No invariance under arbitrary reordered tuples is being presumed.

Define the $G$-invariant cochain $b=cH_{s-1}$. Equation (DC.3) and
$dc=0$ give
\[
 db=c-(\bigwedge h)^*c.
\]
For vertices in $F(Q)$, (DC.1) implies that their images under $h$
lie in the vector space spanned by $j(Q)$, which has dimension
$|Q|$. Thus the second term vanishes on degree-$s$ chains of
$F(Q)$ whenever $s\geq |Q|$. Consequently $b|_{F(Q)}$ is a
primitive of its cocycle. The same holds for $F(P)$, since
$|P|\leq |Q|$. Every $a_i$ extends to an element of $G$ by
homogeneity. The $G$-invariance of $b$ therefore gives
\[
 F(a_i)^*(b|_{F(Q)})=b|_{F(P)}
 \qquad(0\leq i\leq k).
\]
Proposition~\ref{C4:rational-primitive-criterion} now supplies the
integral affine equalizer.

Finally every finite parallel family in a filtered colimit of
categories can be represented at one finite stage: its finitely
many objects, arrows, and required equalities all occur at some
common stage. Represent a family at stage $n$ with target $Q$,
and choose $m\geq n$ with $r_m\geq |Q|$. Apply the first assertion
to the composite functor from stage $n$ to stage $m$ and then
pass its integral equalizer to the colimit. Common successors and
nonemptiness also pass to this colimit. The all-family affine
criterion proves integral exactness. This proof does not assume
that the transition functors are faithful or injective on objects.
\end{proof}



% C.24, strengthened in place; coefficients are genuine Q-vector spaces.
\begin{theorem}[Proved: rational-vector coefficient cocycle colimits]
\label{C4:acyclic-cocycle-colimits}\label{C4:vector-cocycle-colimits}
Let $I$ be a nonempty small filtered category. For each $i\in I$,
let $V_i$ be a countable rational vector space and let $r_i\geq2$.
Allow arbitrary functors in a diagram
$i\mapsto\mathcal A_{r_i}(V_i)$, without requiring that they come
from linear maps of the coefficient spaces. For its categorical
colimit $A$, one has
\[
 \widetilde H_*(BA;\mathbb Q)=0
 \quad\Longleftrightarrow\quad
 A\text{ is exact over }\mathbb Z.
\]
Thus every total, or merely weakly contractible-base, colimit of
this form is integrally exact, including when the coefficient vector
spaces have unbounded or infinite dimension.
\end{theorem}
\begin{proof}
We give the finite obstruction and its functorial control, which
make the argument work in arbitrary coefficient dimension.
For a finite family $f=(a_0,\ldots,a_k):P\to Q$ in
$\mathcal A_r(V)$, let $K_f^*$ be the subcomplex
\[
 K_f^n=\ker\left(
 C^n(\Delta^Q;\mathbb Q)\longrightarrow
 \bigoplus_{i=1}^k C^n(\Delta^P;\mathbb Q)\right),
 \qquad b\longmapsto(a_i^*b-a_0^*b)_{i=1}^k.
 \tag{VC.1}
\]
The prescribed cocycle $c_Q$ belongs to $K_f^r\otimes V$.
Its class
\[
 o_f=[c_Q]\in H^r(K_f^*)\otimes_{\mathbb Q}V
 \tag{VC.2}
\]
vanishes precisely when there is a primitive $b_Q$ with all
$a_i^*b_Q$ equal. Put $b_P=a_0^*b_Q$ in that case. The rational
primitive criterion of Proposition~\ref{C4:rational-primitive-criterion}
works verbatim for $V$: rational affine barycenters of $V$-valued
primitives are defined, and the expansion by arbitrary
$V$-valued $(r-1)$-cochains is the same ordered free-amalgamation
Ramsey class. Hence
\[
 o_f=0
 \quad\Longleftrightarrow\quad
 f\text{ has an integral affine equalizer}.
 \tag{VC.3}
\]
The vector space $H^r(K_f^*)$ is finite dimensional, since $P,Q$
are finite. Consequently (VC.2) involves only finitely many
vectors of $V$, even if $V$ is infinite dimensional.

We record explicitly the cochain comparison, including its
naturality for arbitrary functors. For
$F:\mathcal A_r(V)\to\mathcal A_s(W)$ put
$K_q^F(B)=C_q(\Delta^{\operatorname{vert}F(B)};\mathbb Q)$ and
form the first-quadrant double complex
\[
 \mathcal B_{p,q}(F)=
 \bigoplus_{B_0\to\cdots\to B_p}K_q^F(B_0).
 \tag{VC.4}
\]
Its vertical augmentation is a quasi-isomorphism to the
constant-coefficient bar complex of $\mathcal A_r(V)$, because
each full simplex is nonempty and contractible. The global
target cocycle defines a cochain in bidegree $(0,s)$ of
$\operatorname{Hom}_{\mathbb Q}(\mathcal B(F),W)$; its horizontal
coboundary is zero by naturality. There is a map to the
corresponding double complex for the identity functor of the
target: send each bar string to its $F$-image and retain its
simplex. On vertical augmentations this map is the ordinary
bar map of $F$. In the target identity double complex, every
tuple action is a disjoint union of representables, and the
horizontal degree-zero quotient is the semisimplicial complex
of $K(W,s)$. Its tautological cocycle is the universal
$W$-valued class: on the difference of top simplices with labels
$w$ and $0$ it evaluates to $w$, and their boundaries agree.
Consequently the class of the global cocycle in (VC.4) is
exactly the pullback under $BF$ of this universal class.

If the smooth vertex set $S$ of $F$ has no fixed vertex, every
nonempty increasing tuple in $S$ has a proper open stabilizer.
Theorem~\ref{B:saturation} and
Proposition~\ref{B:cocycle-audit} therefore make every component
of its category of elements contractible. Thus the positive
horizontal homology of (VC.4) vanishes. Its degree-zero
horizontal quotient is $C_*(\Delta^S;\mathbb Q)_G$: every finite
tuple occurs at a finite source stage, and homogeneity realizes
the identifications. Projection to that quotient is a
quasi-isomorphism. Dualizing into $W$ identifies its invariant
cochain cohomology with $H^*(B\mathcal A_r(V);W)$, and the
preceding bar map identifies the cocycle class naturally with
the pulled-back universal class. All total-complex diagonals
are finite, and duality into a rational vector space is exact,
so these comparisons require no convergence or finite-dimension
assumption on $V,W$. If $S$ has a fixed vertex, invariant coning
instead gives a primitive already in bidegree $(0,s-1)$ of
the double cochain complex, so the pullback class is zero.

Consider now any functor
$F:\mathcal A_r(V)\to\mathcal A_r(W)$, and write
\[
 L=(BF)_*:V=H_r(B\mathcal A_r(V);\mathbb Q)
       \longrightarrow H_r(B\mathcal A_r(W);\mathbb Q)=W.
 \tag{VC.5}
\]
We claim that
\[
 F(f)\text{ has an integral affine equalizer}
 \quad\Longleftrightarrow\quad
 (1\otimes L)o_f=0.
 \tag{VC.6}
\]
Use the smooth ordered vertex set $S$ of $F$, with its invariant
$W$-valued cocycle $c_S$, as in Theorem~\ref{C4:degree-colimit-exclusion}.
If $S$ has a fixed vertex, invariant coning makes $c_S$ exact,
and every image family equalizes. The natural bar/simplex
comparison also gives $L=0$, so (VC.6) holds in this case.

Suppose that $S$ has no fixed vertex. The finite-support argument
from Theorem~\ref{C4:degree-colimit-exclusion} applies without
change to $\mathcal A_r(V)$. It supplies an equivariant point
section $j:M\to S$ and the augmentation-preserving support map
\[
 p:\mathbb Q[S]\longrightarrow\mathbb Q[M],
 \qquad p(v)=|B_v|^{-1}\sum_{x\in B_v}x,
 \qquad pj=\operatorname{id}.
\]
For every finite source object $Q$, the support of every vertex
of $F(Q)$ is contained in $Q$. Thus $p$ restricts to a chain map
from the full simplex on $F(Q)$ to that on $Q$, and these maps
are natural under every source embedding. The section $j$ also
restricts naturally: choose one vertex of $F$ of the one-point
object and transport it along its unique embeddings.
It need not preserve the displayed linear orders; the induced
map of alternating simplex chains is nevertheless defined and
natural.

The ordered prism gives an equivariant homotopy $H$ from
$(jp)_*$ to the identity on the simplex chain complex of $S$.
The invariant cochain $j^*c_S$ represents the homomorphism $L$
in
\[
 H^r\bigl(C^*(\Delta^M;W)^G\bigr)
   =H^r(B\mathcal A_r(V);W)
   =\operatorname{Hom}_{\mathbb Q}(V,W).
 \tag{VC.7}
\]
Here the first identification is obtained by the bar/simplex
comparison above for the identity source vertex functor. The second is the
universal coefficient theorem over $\mathbb Q$, since the
classifying space is $K(V,r)$. Naturality follows from the double-complex maps and the
objectwise natural section $j$, whose vertical augmentation
induces the identity on the source bar complex. In particular, for the tautological
$V$-valued source cocycle $c_M$, there is an invariant cochain
$\alpha$ such that
\[
 j^*c_S=Lc_M+d\alpha.
\]
The prism identity therefore gives the global equality
\[
 c_S=p^*(Lc_M)+d\bigl(c_SH+p^*\alpha\bigr).
 \tag{VC.8}
\]
The error primitive in (VC.8) is invariant, so its restrictions
to $F(P),F(Q)$ commute with every arrow in the given family.
If $(1\otimes L)o_f=0$, choose compatible primitives of $Lc_Q$
and pull them back by $p$, adding the error primitive from
(VC.8). This supplies compatible primitives for $F(f)$.
Conversely, pull any compatible primitive for $F(f)$ back by
$j$ and subtract $\alpha$. This supplies a compatible primitive
of $Lc_Q$. This proves (VC.6) by (VC.3).

Two further consequences of the same construction will be used. A functor
$\mathcal A_r(V)\to\mathcal A_s(W)$ with $s\geq|Q|$ sends a
family with target $Q$ to an integrally equalizable family,
by the support-prism rank bound. If $s<r$, the invariant target
cocycle is exact because
$H^s(K(V,r);W)=0$; hence every image family equalizes.
If $V=0$, the source category already has the Ramsey property,
so every image family has an integral affine equalizer. If $W=0$,
the target has that property. Neither case can occur along a
persistent nonequalizable family.

Assume that $BA$ is rationally acyclic, but that $A$ is not
integrally exact. A nonequalizable finite family is represented
at some stage $i$ and remains nonequalizable at every later
stage. Its fixed finite target bounds all future degrees by the
rank bound. Choose a stage $j$ attaining the largest degree in
this filtered tail. No later degree can be smaller, by the
preceding vanishing assertion; hence all degrees on the
cofinal category $j/I$ equal some $r$.

On this tail, filtered categorical colimits commute with nerves
and rational homology, so
\[
 0=H_r(BA;\mathbb Q)=\operatorname*{colim}_{k\in j/I} V_k.
\]
Write the obstruction $o_{f_j}$ in a basis of the finite
dimensional space $H^r(K_{f_j}^*)$. Each of its finitely many
coordinates in $V_j$ becomes zero at some stage. Filteredness
provides one later stage where they all vanish. Apply (VC.6) to
the transition functor from $j$ to that stage: the image family
has an integral affine equalizer, a contradiction.
The converse follows from integral exactness applied to the
constant rational coefficient diagram.
\end{proof}


\subsection{Filtered totalizations of Kan simplex categories}

\begin{theorem}[Proved: contractible filtered Kan totalizations are exact]
\label{C:kan-filtered-totalization}
Let $I$ be a nonempty small filtered category, let $K_i$ be nonempty
connected small Kan complexes, and let
\[
 A_i=\Delta_{\mathrm{inj}}/K_i\quad(i\in I)
\]
form a diagram of categories with arbitrary transition functors.
The transition functors need not be faithful or induced by simplicial
maps. Put $A=\colim_{i\in I}A_i$, the colimit in small categories.
Then the following conditions are equivalent:
\[
 BA\simeq *;\qquad A\text{ is total};\qquad
 \colim_A\text{ is exact over }\Z.
\]
When these conditions hold, $\colim_A$ is exact over every unital ring.
Consequently no total counterexample to C.4, and no total nonexact
category, can be obtained as such a filtered colimit. This includes
arbitrary coefficient groups and arbitrary connected Kan stages;
it imposes no increasing-connectivity condition on the stages.
\end{theorem}
\begin{proof}
We prove the nontrivial implication from weak contractibility of $BA$
to integral exactness. We first give the finite obstruction and its
functorial comparison, including the distinction between ordinary
and fat realization.

\emph{A finite family and a contractible path-space lift.}
Write $A(K)=\Delta_{\mathrm{inj}}/K$. For a finite parallel family
$f=(a_0,\ldots,a_m):P\rightrightarrows Q$ in $A(K)$, put
\[
 R_f=\operatorname{coeq}\left(
 \coprod_{j=1}^{m}\Delta^{P}\;
 \substack{\longrightarrow\\[-5pt]\longrightarrow}\;
 \Delta^{Q}\right).
 \tag{KC.1}
\]
Here $\Delta^B$ is the standard simplicial simplex on the finite
ordered vertex set of $B$, and the two arrows on the $j$th summand
are $a_j$ and $a_0$. The coequalizer is taken in simplicial sets.
It has finitely many nondegenerate simplices, although vertices or
faces can be identified. Its realization is therefore a finite CW
complex. All generating identifications are induced by injective
face maps and preserve degeneracy type, so a nondegenerate cell is
not identified with a degenerate simplex. The simplex label of $Q$ induces a simplicial map
$c_f:R_f\to K$.

If $|c_f|$ is nullhomotopic, the family $f$ has an integral affine
equalizer. Indeed, choose a vertex $k\in K$ and the based path-space
Kan fibration $\pi:P_kK\to K$. Its total space is weakly contractible;
the precise construction and proof are
\cite[Lemma I.7.5, pp.~30--31]{ABGoerssJardine}. The map
\[
 \operatorname{Map}(R_f,P_kK)\longrightarrow
 \operatorname{Map}(R_f,K)
\]
is a Kan fibration by \cite[Proposition I.5.2 and Corollary I.5.3]{ABGoerssJardine}.
Its image on vertices is a union of path components: lift an edge
starting at a lifted vertex, and repeat along an edge path. The
constant map at $k$ lifts, and $c_f$ belongs to its component.
Thus $c_f$ has a strict simplicial lift to $P_kK$. Pulling this lift
back to $\Delta^P$ and $\Delta^Q$ gives objects $P^*,Q^*$ and the
entire parallel family in $A(P_kK)$. This category is Ramsey and
integrally exact by Theorem~\ref{C4:contractible-kan-ramsey}.
An integral affine equalizer there is carried to one in $A(K)$ by
$\pi$. All augmentations are preserved. We need only this sufficient
nullhomotopy criterion; we do not assert its converse for arbitrary
$K$.

\emph{The map of base spaces associated with an arbitrary functor.}
Let $F:A(K)\to A(L)$, where $K$ is simply connected and $L$ is
connected. Denote by $\|K\|$ the fat realization retaining degenerate
simplices as cells. Then
\[
 \|K\|=\colim_{B\in A(K)}|\Delta^B|,
 \qquad BA(K)\simeq\|K\|\simeq |K|.
 \tag{KC.2}
\]
For clarity, the first weak equivalence can be obtained from the
homotopy colimit of the displayed simplex diagram. Its augmentation
to $BA(K)$ is a weak equivalence: in each bar degree it is the
disjoint union, over strings of arrows, of the weak equivalences
from a simplex to a point. Its map to the
ordinary colimit is a weak equivalence: order the index by simplex
dimension; at each object the colimit of its proper faces is exactly
the boundary of that simplex, so the diagram is built by cofibrations.
The second equivalence is the degeneracy comparison already checked
in Theorem~\ref{B:kan-recovery}. This also specifies the comparison
used below; no invariance under an arbitrary functor is presumed.

For every vertex label $v\in K_0$, choose a vertex $z_v$ of the
nonempty object $F(\{v\})$. For every object $B$, the assignments
\[
 j_B(x)=F(\{x\}\hookrightarrow B)(z_{\operatorname{label}(x)})
 \quad(x\in B)
\]
extend affinely to continuous maps
$j_B:|\Delta^B|\to|\Delta^{F B}|$. They are natural for all arrows
of $A(K)$. They need not preserve vertex orders, which causes no
problem for maps of geometric simplices. Composing with the target
simplex labels gives
\[
 J_F:\|K\|\longrightarrow\|L\|.
 \tag{KC.3}
\]
The same natural maps on homotopy colimits, followed by the change of
index along $F$, lie over $BF$ on the base augmentations. Hence
$J_F$ represents $BF$ under (KC.2). The canonical map
$\bar c_f:|R_f|\to\|K\|$ is obtained from the cell of $Q$;
its composite to $|K|$ is $|c_f|$.

We claim that if $J_F\bar c_f$ is nullhomotopic, then $Ff$ has an
integral affine equalizer. First suppose $K$ is countable and use
its homogeneous limit $M$ and automorphism group $G$. Theorem
\ref{B:kan-recovery} gives strong amalgamation, (St), (Op), and
$\operatorname{Fix}_M(G_B)=B$ for every finite $B\subset M$.
Every finite subset is an admissible substructure. Form the smooth
ordered $G$-set
\[
 S=\colim_{B\subset M,\ B\ne\varnothing}
       V(FB),
\]
where this colimit uses only inclusions. The maps are injective, so
we regard the vertex sets as subsets of $S$. Every finite subset of
$S$ lies in some $V(FB)$. Their coherent labels give a continuous
$G$-invariant map $c_S:|\Delta^S|\to\|L\|$.

If $S$ has a $G$-fixed vertex $z$, cone each point of
$|\Delta^{FQ}|$ linearly to $z$ in $|\Delta^S|$, then apply $c_S$.
This homotopy descends to $|R_{Ff}|$. Indeed, extend each
$a_j a_0^{-1}$ to an element $g_j\in G$; it sends the image of
$F a_0$ to that of $F a_j$, fixes $z$, and leaves $c_S$ invariant.
The resulting homotopy is a nullhomotopy of $\bar c_{Ff}$, so the
path-space criterion applies.

Suppose now that $S$ has no fixed vertex. By (Op), every $s\in S$
has a unique nonempty finite support $B_s$ with
$G_s=G_{B_s}$. If $s\in V(FB)$ then $G_B\leq G_s$, and
$\operatorname{Fix}_M(G_B)=B$ implies $B_s\subset B$.
The natural maps $j_B$ above assemble to a $G$-equivariant vertex
map $j:M\to S$. A point $j(x)$ has support exactly $\{x\}$:
its support is contained in that singleton and cannot be empty.
In particular $j$ is injective. Define the positive barycentric map
\[
 p:|\Delta^S|\longrightarrow|\Delta^M|,
 \qquad p(s)=\frac1{|B_s|}\sum_{x\in B_s}x.
 \tag{KC.4}
\]
This is a continuous affine map on every finite simplex, hence
continuous for the CW topology. It is $G$-equivariant and restricts
to $p_B:|\Delta^{FB}|\to|\Delta^B|$. These restrictions are natural
under every embedding, by extension to $G$. We have $p_Bj_B=1$,
and the straight-line homotopy
\[
 H_B(u,t)=(1-t)u+t\,j_Bp_B(u)
 \tag{KC.5}
\]
stays in $|\Delta^{FB}|$ and is natural under all arrows.
Consequently $p_B,j_B,H_B$ descend for the family to maps
$p_f:|R_{Ff}|\to|R_f|$, $j_f:|R_f|\to|R_{Ff}|$ and a homotopy
$j_fp_f\simeq1$. The label maps satisfy
\[
 \bar c_{Ff}\simeq J_F\bar c_f p_f.
 \tag{KC.6}
\]
Therefore the assumed nullhomotopy of $J_F\bar c_f$ gives a
nullhomotopy of $c_{Ff}$ after passage to ordinary realization,
and hence an integral affine equalizer.

The countability restriction on $K$ is unnecessary for this local
claim. Any specified finite family lies in $A(K')$ for a countable
simply connected Kan subcomplex $K'\subset K$: start with its
finitely many labels, close under degeneracies and faces, fill the
countably many horns, add paths and nullhomotopies for all edge
loops, and repeat. The ambient connected simply connected $K$
supplies these choices, and their countable union is Kan and simply
connected. Apply the preceding argument to $F|_{A(K')}$; its image
vertex set is countable regardless of the size of $L$. The finite
family, the chosen $j_B$, and its nullhomotopy in $\|L\|$ are
unchanged. This proves the claim for all small simply connected $K$.

\emph{Passage to a filtered stage.}
For a filtered diagram of small categories, the natural map
$\colim_i N A_i\to N(\colim_i A_i)$ is an isomorphism: each
composable finite string and each equality between two such strings
is represented at one stage. Homotopy classes from a finite
simplicial set commute with these filtered colimits after fibrant
replacement. One explicit verification uses $\operatorname{Ex}^\infty$:
it is a natural Kan replacement by
\cite[Theorem III.4.8, p.~188]{ABGoerssJardine}, and it commutes
with filtered colimits because each $\operatorname{Sd}\Delta^n$
is finite. A map from a finite simplicial set, and a homotopy from
it, therefore use only finitely many simplices and equalities and
are realized at one stage. This applies also to the finite sets
$R_f$ and to a simplicial circle.

Assume $BA\simeq *$ and take a finite parallel family in $A$.
It lifts to a parallel family $f:P\rightrightarrows Q$ in some
$A_i$. The finitely many loops $a_j a_0^{-1}$ in the fundamental
groupoid of $BA_i$ become trivial at one later stage $A_j$.
Replace $f$ by its image there. The universal simplicial cover
$\widetilde K_j\to K_j$ induces the universal covering category
$A(\widetilde K_j)\to A(K_j)$, as checked in
Theorem~\ref{B:kan-recovery}. Choose a lift of $Q$. Pulling it back
along $a_0$ chooses a lift of $P$; the trivial loop classes say that
pullback along every $a_j$ gives this same lift. Thus the whole
parallel family lifts to a family $\widetilde f$ in
$A(\widetilde K_j)$.

The finite complex $R_{\widetilde f}$ maps by its fat label into
$BA(\widetilde K_j)$, using (KC.2), and then into $BA_j$.
Its image in the weakly contractible $BA$ is nullhomotopic.
The finite-domain argument gives a further transition $j\to k$
for which this composite is already nullhomotopic in $BA_k$.
Apply the local claim to the functor
\[
 A(\widetilde K_j)\longrightarrow A(K_j)
       \longrightarrow A(K_k).
\]
Its source Kan complex is simply connected, its target is connected,
and its base map kills the indicated finite label. The image of
$\widetilde f$, which is the image of the original family, has an
integral affine equalizer. It remains an affine equalizer in $A$.

Every $A_i$ has common successors, by connectedness and the Kan
amalgamation proof in Theorem~\ref{B:kan-recovery}; this property
passes to the filtered colimit. We have proved the integral affine
equalizer criterion for every finite parallel family, so $\colim_A$
is integrally exact. The same integer witnesses give exactness over
every unital ring, including the zero ring. Integral exactness
implies totality by Theorem~\ref{F:companion-chain}, and totality
implies $BA\simeq *$ by applying its definition to the constant
singleton action. This proves all the equivalences.
\end{proof}


% Independent exclusion of nonfiltered coefficient-pushforward assemblies.
\begin{theorem}[Proved: cocycle assemblies over an exact base]
\label{C4:grothendieck-exclusion}
Let $B$ be a nonempty connected small category that is exact over
$\mathbb Z$, let $\Lambda:B\to\Ab$ be a diagram of small abelian
groups, and fix $r\geq2$.
Form the covariant Grothendieck construction
\[
 A=\int_{b\in B}\mathcal A_r(\Lambda(b)),
 \tag{GA.1}
\]
using the coefficient-pushforward functors associated with the
homomorphisms of $\Lambda$. The following conditions are equivalent:
\begin{enumerate}
\item $\Lambda$ is locally zero: every $v\in \Lambda(b)$ is killed by
some arrow out of $b$;
\item $A$ is exact over $\mathbb Z$;
\item $A$ is total.
\end{enumerate}
If every $\Lambda(b)$ is uniquely $p$-divisible, then $A$ is
exact over $\mathbb F_p$. In particular, rational-vector coefficient
assemblies are exact over every prime field, but they give no total
counterexample to the all-prime implication.
\end{theorem}
\begin{proof}
The objects of $A$ are pairs $(b,P)$ with
$P\in\mathcal A_r(\Lambda(b))$. An arrow $(b,P)\to(c,Q)$ is a
pair $(a,\alpha)$, where $a:b\to c$ and $\alpha$ is an
increasing vertex injection whose restriction of the target
cocycle is $\Lambda(a)c_P$. This description fixes the variance in
(GA.1).

Under the stated uniquely $p$-divisible hypothesis, each fibre
is exact over $\mathbb F_p$, by the perfect-group-algebra proof.
For an uncountable coefficient group, any finite family lies in
the countable subgroup generated by its labels and all their
iterated $p$th roots, so the same finite-family criterion applies.
For any module
diagram $N$ on $A$, its colimit is the colimit over $B$ of the
fibrewise colimits. This equality follows directly from cocones:
each base arrow acts on a fibre colimit through the natural
cocartesian arrows in the Grothendieck construction. Each
fibrewise colimit is exact over $\mathbb F_p$, and the colimit
over $B$ is exact over every ring by its integral affine
equalizers. Their composite is exact. This proves the
prime-field assertion for all module diagrams.

Suppose first that $\Lambda$ is locally zero. A finite set of elements
at one object is killed by one common arrow. Indeed exactness
of $B$ implies condition I, and an amalgamating square for two
arrows that kill specified elements gives one composite killing
both; induction treats a finite set.
Take a finite parallel family $(a_i,\alpha_i):(b,P)\to(c,Q)$.
Choose $h:c\to d$ killing every one of the finitely many
cocycle labels on $Q$, and postcompose the family with
$(h,\operatorname{id}_Q):(c,Q)\to(d,Q_0)$, where $Q_0$
has zero cocycle. Consider the same arrow pairs
$(ha_i,\alpha_i)$ as a family from $(b,P_0)$ to $(d,Q_0)$
in the zero-cocycle subcategory, which is $B\times\Delta_{\rm inj}$.
This product is integrally exact: its module colimit is the
successive colimit over the two exact factors. Its finite-family
affine criterion therefore gives an integral affine equalizer
of these arrow pairs. The equalizer's outgoing arrows have
zero-cocycle source and target and so are also arrows of $A$.
Equality of their composites is equality of the base arrow
and the increasing vertex map, independently of the original
source cocycle. Hence the same affine combination equalizes
the original postcomposed family, and then the original family.
The common-successor property for $A$ follows by moving base
objects to a common successor and then using joint embedding
in that fibre. Thus the affine criterion proves integral
exactness of $A$.

We show that a failure of local zero contradicts totality.
Choose $v$ that no outgoing arrow kills. If some reachable image
has finite order, choose the least such order $n\geq2$ and an
image of order $n$. Every further image still has order $n$:
its order divides $n$, cannot be one, and cannot be smaller by
minimality. If no image has finite order, put $n=\infty$.
Define a covariant set action by
\[
 L(b)=\{w\in\Lambda(b):\operatorname{ord}(\Lambda(a)w)=n
             \text{ for every arrow }a\text{ out of }b\}.
 \tag{GA.2}
\]
It is nonempty somewhere and functorial by composition. Put
$C_n=\mathbb Z/n$ for finite $n$ and $C_\infty=\mathbb Z$.
On $\operatorname{El}_B(L)$ the pulled-back coefficient diagram
has a tautological injective natural map from the constant
$C_n$, given by $1\mapsto w$ at $(b,w)$.

Every category of elements of a set action on an exact category
is exact. To check this explicitly, for its projection $u$ the
left Kan extension of a module diagram is
\[
 (u_!N)(b)=\bigoplus_{v\in L(b)}N(b,v).
\]
The summand formula holds because each component of the comma
category $u/b$ has the terminal object determined by its image
element in $L(b)$. Thus $u_!$ is exact, and
$\operatorname{colim}_{\operatorname{El}_B L}=
\operatorname{colim}_B\,u_!$ is exact. Each connected component
$D$ is exact as well. Choose a nonempty such $D$; its constant
coefficient colimit is $C_n\ne0$, so exactness and the
tautological injection give
\[
 0\ne C_n\lhook\joinrel\longrightarrow
    \operatorname{colim}_{(b,v)\in D}\Lambda(b).
 \tag{GA.3}
\]

Pull the action $L$ back along $A\to B$. The part of its
category of elements over $D$ is exactly
\[
 E=\int_{(b,v)\in D}\mathcal A_r(\Lambda(b)).
\]
It is connected, since $D$ and all fibres are connected and
each fibre is nonempty. The Grothendieck-construction homology
spectral sequence, or the bar double complex for the fibre
nerves, gives
\[
 E^2_{p,q}=L_p\operatorname{colim}_{D}
 H_q(B\mathcal A_r(\Lambda(b));\mathbb Z)
 \ \Longrightarrow\ H_{p+q}(BE;\mathbb Z).
 \tag{GA.4}
\]
Here is a direct justification of this comparison. For the projection
$\pi:E\to D$, the comma category $\pi/d$ retracts to its fibre:
send $(b,X,a:b\to d)$ to the coefficient pushforward $a_*X$,
and include the fibre using $a=1_d$. The cocartesian arrow
$(a,1)$ gives a natural transformation from the identity to the
inclusion followed by this retraction. Thus the bar formula for
derived left Kan extension gives
$L_q\pi_!\mathbb Z(d)=H_q(B\mathcal A_r(\Lambda(d));\mathbb Z)$,
naturally in $d$. Resolve these bar complexes by representable
modules over $D$ and compute the resulting first-quadrant double
complex in its two orders. Since $\pi_!$ preserves projectives
(its right adjoint, restriction, is exact), the total complex computes
$L_*\operatorname{colim}_E\mathbb Z$. This gives exactly (GA.4).
Because $D$ is exact, all columns $p>0$ vanish. Since
$B\mathcal A_r(\Lambda(b))\simeq K(\Lambda(b),r)$ naturally for
coefficient-pushforward maps, (GA.3)--(GA.4) imply
\[
 H_r(BE;\mathbb Z)=
 \operatorname{colim}_{(b,v)\in D}\Lambda(b)\ne0.
\]
This connected action component is not weakly contractible,
so $A$ is not total. Finally integral exactness implies
totality, completing the equivalences.
\end{proof}


% Independently audited and integrated as C.27.
\begin{theorem}[Proved: Kan assemblies over an exact base]
\label{C4:exact-base-kan-assembly}
Let $B$ be a nonempty connected small category exact over
$\mathbb Z$. Let $b\mapsto C_b=A(K_b)=\Delta_{\rm inj}/K_b$
be a diagram of categories and arbitrary functors, where every
$K_b$ is a nonempty connected small Kan complex. Put
$A=\int_B C_b$. The following conditions are equivalent:
\begin{enumerate}
\item $A$ is total;
\item $A$ is exact over $\mathbb Z$;
\item every positive homology diagram
$b\mapsto H_q(BC_b;\mathbb Z)$ is locally zero, and every
loop in every $BC_b$ becomes trivial after some base arrow.
\end{enumerate}
In the last condition, a base arrow acts by the specified
transition functor; the target basepoint is its image of the
given basepoint. In particular, allowing arbitrary nonfiltered
exact bases and arbitrary transition functors still produces
no total nonexact category by such an assembly.
\end{theorem}
\begin{proof}
Write $F_u:C_b\to C_c$ for the functor associated with
$u:b\to c$. The transition diagram is an ordinary strict
diagram; the Grothendieck construction has arrows
$(u,\alpha):(b,x)\to(c,y)$ with
$\alpha:F_u x\to y$.

\emph{The derived fibre formula, including nonconstant coefficients.}
For any diagram $N:A\to\Ab$ and the projection $\pi:A\to B$,
one has natural isomorphisms
\[
 (L_q\pi_!N)(b)=T_q^{C_b}(N|_{C_b}),\qquad
 T_q^A(N)=\operatorname{colim}_{b\in B}
             T_q^{C_b}(N|_{C_b}).
 \tag{EA.1}
\]
We spell out the first assertion to control coefficients.
An object of $\pi/b$ is $(c,x,u:c\to b)$. The fibre inclusion
$i:C_b\to\pi/b$ has left adjoint
\[
 R(c,x,u)=F_u x.
\]
The unit is the cocartesian arrow
$(c,x,u)\to(b,F_u x,\operatorname{id}_b)$, and $Ri=1$.
On bar complexes with the pulled-back coefficients, the map
induced by $R$ first applies $N$ to this unit at the initial
vertex of a bar string. Inclusion induces the reverse map.
One composite is the identity, and the other is chain homotopic
to the identity by the usual natural-transformation prism for
the unit. Thus these are coefficient-preserving derived
colimit comparisons, not only a constant-coefficient homology
calculation.

The pointwise derived left-Kan formula can be computed with
the same comma bar complex. For completeness, the restriction
to $\pi/b$ of a free representable $\mathbb Z[A(a_0,-)]$ is
colimit-acyclic: its category of elements splits according to
the composite arrow $\pi(a_0)\to b$, and each part has the
initial object given by $a_0$, that base arrow, and
$\operatorname{id}_{a_0}$. Projective resolutions therefore
justify the pointwise formula for arbitrary $N$. Finally
$\operatorname{colim}_A=\operatorname{colim}_B\pi_!$.
Since the second colimit is exact, it commutes with homology
of a projective resolution, giving the second equality in
(EA.1).

\emph{Persistent homology contradicts totality.}
Every category of elements of a set action on $B$ is exact:
its left Kan extension to $B$ is the exact fibrewise direct
sum, as checked in Theorem~\ref{C4:grothendieck-exclusion}.
Let $H(b)=H_q(BC_b;\mathbb Z)$ for $q>0$. If an element is
never killed by a base arrow, move it, if necessary, to have
the least finite order attainable by an image; if no finite
order is attainable, all its images have infinite order.
Thus for some $n\in\{2,3,\ldots,\infty\}$ the set action
\[
 T(b)=\{v\in H(b):\operatorname{ord}(H(u)v)=n
                    \text{ for every }u\text{ out of }b\}
\]
is nonempty somewhere. On a nonempty connected component
$D$ of $\operatorname{El}_B T$, the pulled-back homology
diagram contains the constant cyclic group $C_n$ injectively,
via $1\mapsto v$; here $C_\infty=\mathbb Z$. Exactness of
$D$ gives a nonzero injection of $C_n$ into its colimit.
Pull $T$ back along $A\to B$. The action component over $D$
is $\int_D C_b$, and is connected because the base and fibres
are connected. Formula (EA.1), now over $D$ and with constant
coefficients, shows that its $q$th homology is
$\operatorname{colim}_D H(b)\ne0$. Totality is contradicted.
This proves the homological part of (iii).

\emph{Persistent fundamental-group loops contradict totality.}
Let $\Pi_b$ be the fundamental groupoid of $BC_b$, on the
vertices that are objects of $C_b$. The functors $F_u$ induce
strict functors of these groupoids. Suppose that a loop
$g\in\Pi_b(z,z)$ is never sent to the identity by a base arrow.
As in the preceding paragraph, move to a loop whose order
$n$ persists under every further base arrow. Define
\[
 T(b)=\{(z,g): z\in C_b,\ g\in\Pi_b(z,z),
       \ \operatorname{ord}(F_u g)=n
          \text{ for every }u\text{ out of }b\}.
 \tag{EA.2}
\]
This is a nonempty covariant set action on $B$.
There is a covariant set action $L$ on the whole assembly,
defined at $(b,x)$ by
\[
 L(b,x)=\coprod_{(z,g)\in T(b)}
      \Pi_b(z,x)/\langle g\rangle.
 \tag{EA.3}
\]
The quotient in (EA.3) is by right composition: $h$ and
$h\circ g^k$ represent the same element. For
$(u,\alpha):(b,x)\to(c,y)$ the action sends
\[
 (z,g,[h])\longmapsto
       (F_u z,F_u g,[\alpha\circ F_u h]).
 \tag{EA.4}
\]
Right powers map to right powers, so (EA.4) is well defined;
strict composition of the transition functors and left
composition by fibre arrows verify the action law.

For a fixed $(z,g)$ the corresponding component of the fibre
action category is the connected cover of $BC_b$ with
fundamental group $\langle g\rangle$. Indeed the fibre
action is the usual path-groupoid coset action, based at
$(z,[\operatorname{id}_z])$; all fibre arrows act bijectively.
Its first homology is canonically $C_n$, with $1$ represented
by $g$. A base arrow carries this chosen basepoint and loop
to the chosen basepoint and loop for $(F_u z,F_u g)$.
Since their orders remain $n$, the map on these first
homology groups is the identity under the stated identifications.

Choose a connected component $D$ of $\operatorname{El}_B T$.
The corresponding part of $\operatorname{El}_A L$ is a
Grothendieck construction over the exact category $D$, with
these connected covering categories as fibres. It is connected.
Formula (EA.1) gives its first homology as
$\operatorname{colim}_D C_n=C_n\ne0$. This contradicts
totality. Therefore every fibre loop is killed by some base
arrow, completing (i)$\Rightarrow$(iii).

\emph{A sequence that kills a given derived-colimit cycle.}
Assume (iii), and fix $N:A\to\Ab$, a base object $b_0$, and
a positive-degree cycle in the bar complex of
$(C_{b_0},N|_{C_{b_0}})$. We construct a sequence of base
arrows and countable connected Kan subcomplexes
\[
 b_0\longrightarrow b_1\longrightarrow b_2\longrightarrow\cdots,
 \qquad K'_j\subset K_{b_j},
 \qquad C'_j=A(K'_j),
 \tag{EA.5}
\]
such that every transition functor restricts to
$C'_j\to C'_{j+1}$, the initial cycle lies in $C'_0$, and
the categorical colimit $C_\infty=\colim_j C'_j$ has weakly
contractible classifying space.

Start with all labels occurring in the finitely many objects
of the given cycle. They lie in a countable connected Kan
subcomplex: close under faces and degeneracies, add paths to
one vertex, choose horn fillers, and repeat. Each $C'_j$ is
countable. Enumerate its positive integral bar cycles and
its fundamental-groupoid loops. At step $j$, schedule the
first $j$ entries in the enumerations for the first $j$
stages. Their images at the current stage are finitely many
homology classes and loops. By (iii), each is killed by a
base arrow. Exactness of $B$ supplies condition I, so finitely
many such arrows can be amalgamated to one common composite
$b_j\to b_{j+1}$ that kills all scheduled classes and loops.

In the full category $C_{b_{j+1}}$, choose finite bar-boundary
witnesses for the killed cycles and finite nerve nullhomotopies
for the killed loops. These use only finitely many object
labels for each scheduled entry. Choose $K'_{j+1}$ containing
those labels and the labels of the images of every object of
$C'_j$. This is still countable; close it under faces,
degeneracies, connecting paths, and horn fillers as before.
The transition functor now restricts as required, and all
scheduled witnesses lie in $C'_{j+1}$. This construction works
even when the full $K_b$ are uncountable and the transition
functors are not induced by simplicial maps.

Every positive homology class and every loop from any stage
eventually occurs in the schedule and becomes zero. Filtered
categorical colimits commute with nerves and with the finite
homotopy and homology tests used here, as checked in
Theorem~\ref{C:kan-filtered-totalization}. Hence $BC_\infty$
is connected, simply connected, and integrally acyclic, so it
is weakly contractible. That theorem now implies that
$C_\infty$ is integrally exact.

Restrict $N$ to the categories in (EA.5). The cocartesian arrows
of $A$ supply compatible natural maps from each coefficient
diagram to the pullback of the next. They define a coefficient
diagram $N_\infty$ on $C_\infty$: for an object represented
at stage $j$, take the filtered colimit of its groups along
the succeeding images. Two representatives of the same object
agree at a later stage, so the definition is independent of
the representative, and the same argument defines its arrow
maps. Degree by degree,
\[
 \operatorname{colim}_j
    C_*(C'_j;N|_{C'_j})
       \cong C_*(C_\infty;N_\infty).
 \tag{EA.6}
\]
Indeed each bar string, each coefficient, and each equality
involves finitely many data and is represented at one stage.
The original positive-degree cycle is a boundary on the right
of (EA.6), by exactness of $C_\infty$. It is therefore a
boundary at one finite stage on the left. Thus, for every
$q>0$, the diagram $b\mapsto T_q^{C_b}(N|_{C_b})$ is locally
zero.

Formula (EA.1) now gives $T_q^A(N)=0$ for $q>0$, because
the colimit of a locally zero diagram is zero. This proves
(iii)$\Rightarrow$(ii). Integral exactness implies totality,
so (ii)$\Rightarrow$(i) completes the proof.
\end{proof}


% Appended finite obstruction; does not assert an ultraproduct theorem.
\begin{theorem}[Proved: three-arrow obstructions in every cocycle degree]
\label{C4:three-arrow-high-degree}
For every integer $r\geq2$ there are finite objects $P_r,Q_r$ of
$\mathcal A_r(\mathbb Q)$ and three increasing embeddings
$f_0,f_1,f_2:P_r\to Q_r$ which have no rational affine equalizer.
Moreover, their finite simplicial coequalizer
\[
 R_r=\Delta^{Q_r}/(f_0\Delta^{P_r}=f_1\Delta^{P_r}
                         =f_2\Delta^{P_r})
\]
has nonzero rational homology in degree $r$. Thus there is no
bound, depending only on the number of arrows, on the homology
degree of such finite coequalizers or of their compatible-primitive
obstructions. The assertion uses three arrows independently of $r$;
the sizes of the two finite objects are allowed to grow with $r$.
\end{theorem}
\begin{proof}
We give the objects, cocycle, and detecting cycle explicitly. Let
\[
 A=\bigoplus_{i\in\mathbb Z}\mathbb Z e_i,\qquad
 \Gamma=A\rtimes\mathbb Z,
 \qquad(f,m)(g,n)=(f+T^m g,m+n),\quad Te_i=e_{i+1}.
\]
Write $a=(e_0,0)$ and $t=(0,1)$. These two elements generate
$\Gamma$. Give $\Gamma$ its left-invariant total order whose positive
cone consists of $(f,m)$ with $m>0$, or with $m=0$ and the highest
nonzero coordinate of $f$ positive. Indeed the indicated order of
$A$ is invariant under every shift, so this lexicographic extension
is a group order; conjugation also preserves its positive cone.
In particular every $a_i=(e_i,0)=t^iat^{-i}$ is positive. All finite
subsets below carry the induced order.

Define a shift-invariant alternating $r$-linear form $\omega$ on
$\mathbb Q\otimes A$ by setting
\[
 \omega(e_{i_1},\ldots,e_{i_r})=
 \begin{cases}
  \operatorname{sgn}(\pi),&
   (i_{\pi(1)},\ldots,i_{\pi(r)})=(j,j+1,\ldots,j+r-1)
   \text{ for some }j,\\
  0,&\text{otherwise}.
 \end{cases}
 \tag{TA.1}
\]
Here the sign is that of the sorting permutation, and a repeated
index gives zero. This specifies a well-defined form because every
input has finite support. On vertices $g_j=(v_j,m_j)\in\Gamma$ put
\[
 c(g_0,\ldots,g_r)
   =\omega(v_1-v_0,\ldots,v_r-v_0).
 \tag{TA.2}
\]
This is an alternating cochain on the full simplex with vertex set
$\Gamma$. Its coboundary is zero: expansion by multilinearity in
the alternating sum of its $r+2$ faces cancels every term. Equivalently,
the affine volume of the boundary of any $(r+1)$-simplex is zero.
Left multiplication by $(h,k)$ translates all $v_j$ by $h$ and
shifts their differences by $T^k$, so (TA.1) proves that $c$ is
$\Gamma$-invariant.

Let $V$ be the $2^r$ lamp vertices
\[
 V=\left\{\left(\sum_{i=0}^{r-1}\epsilon_i e_i,0\right):
                    \epsilon_i\in\{0,1\}\right\}.
\]
For each $0\leq i<r$, let $V_i^-$ be its face with $\epsilon_i=0$;
the opposite face is $V_i^+=a_i V_i^-$. Realize the word
$a_i=t^iat^{-i}$ as a sequence of left translations: first $i$
steps by $t^{-1}$, then one by $a$, then $i$ by $t$.
Let $P_0$ be the union of $V$ and all intermediate translates of
all the $V_i^-$ along these finitely many paths, including both
endpoints. Put
\[
 Q=P_0\cup aP_0\cup tP_0,\qquad
 P=\{g\in Q:ag\in Q\text{ and }tg\in Q\}.
 \tag{TA.3}
\]
Both are finite, $P_0\subset P\subset Q$, and $P$ is nonempty.
The three increasing injections $f_0(g)=g$, $f_1(g)=ag$,
$f_2(g)=tg$ have equal pullbacks of $c|_Q$, by invariance. Hence
$Q$ with cocycle $c|_Q$, and $P$ with their common pullback, are
objects of $\mathcal A_r(\mathbb Q)$ and these are morphisms.

In the oriented simplex chain group on $V$, use the standard cube
triangulation
\[
 z=\sum_{\pi\in S_r}\operatorname{sgn}(\pi)
 \left[0,e_{\pi(0)},e_{\pi(0)}+e_{\pi(1)},\ldots,
                         e_0+\cdots+e_{r-1}\right].
 \tag{TA.4}
\]
Lamp vectors in this formula mean their vertices in $\Gamma$.
Every displayed tuple is increasing, since its consecutive
differences are positive lamps. Its boundary is the sum of the
oriented triangulations of opposite cube faces, with opposite
signs. Each upper face is identified with its lower face in
$R_r$: along a step by $a$ or $t$ the starting face lies in $P_0$
and is identified with its translate by the corresponding
coequalizer relation; along a step by $t^{-1}$ the ending face
lies in $P_0$, and the same $t$-relation is used in reverse.
The paths specified above therefore identify the whole ordered
face simplices, not merely their vertices. Consequently the image
$\bar z$ of $z$ in the normalized simplicial chains of $R_r$ is a
cycle.

Because the restrictions of $c$ along all three arrows agree,
$c|_Q$ descends to an $r$-cocycle $\bar c$ on $R_r$. Formula
(TA.1) gives
\[
 c\left(0,e_{\pi(0)},e_{\pi(0)}+e_{\pi(1)},\ldots,
                  e_0+\cdots+e_{r-1}\right)
       =\operatorname{sgn}(\pi),
 \qquad \langle\bar c,\bar z\rangle=r!\ne0.
 \tag{TA.5}
\]
Thus $[\bar z]\ne0$ in $H_r(R_r;\mathbb Q)$. In particular the
cocycle on $Q$ has no primitive whose three restrictions to $P$
agree: such a primitive would descend to $R_r$ and make the
pairing in (TA.5) zero. Proposition~\ref{C4:rational-primitive-criterion}
now rules out a rational affine equalizer (and therefore an
integral one).
\end{proof}

\paragraph{Scope of the compactness approach.}
This theorem supplies a fixed number of arrows for rational
nonexactness in arbitrarily high cocycle degree. The ordinary
category ultraproduct nevertheless fails modular exactness:
Theorem~\ref{C4:ultraproduct-exclusion} gives a two-arrow
zero-cocycle family whose affine witnesses over every nonzero
ring require unbounded support. The three-arrow family just
constructed itself inherits this obstruction. Indeed $P_0$
contains all pure shifts $(0,m)$ for $-(r-1)\leq m\leq0$:
use the path for $i=r-1$ starting at the zero vertex of its lower
face. The set $Q$ also contains $(0,1)$. Restricting $f_0,f_2$
to these $r$ vertices gives the first-window and last-window
inclusions into the $r+1$ pure-shift vertices. Their cocycles
are zero. Any affine equalizer of the three arrows therefore
restricts to an affine equalizer of this window pair, with no
increase in support. Its support has size at least $r$, by
\eqref{C4:window-support}. Consequently even a selected subcategory of the
$r\to\infty$ ultraproduct which retains this distinguished
family has no affine equalizer for it over any nonzero ring.
That direct compactness construction cannot give the
total-category counterexample sought in C.4.
The present theorem does not contradict
Theorem~\ref{C4:degree-colimit-exclusion}, whose degree bound
depends on the size of a fixed source target object and whose
transition maps are actual functors.


\begin{theorem}[Proved: ordinary ultraproducts lose modular exactness]
\label{C4:ultraproduct-exclusion}
Let $\mathcal U$ be a nonprincipal ultrafilter on $\mathbb N$.
For arbitrary integers $r_i\geq2$, the ordinary category
ultraproduct
\[
 \mathcal C=\prod_{i\to\mathcal U}\mathcal A_{r_i}(\mathbb Q)
\]
is not exact over any nonzero unital ring and is not total. The same is
true of the ultrapower $\prod_{i\to\mathcal U}\Delta_{\rm inj}$,
although $\Delta_{\rm inj}$ itself is integrally exact and total.
Thus there is no bound on the support size of affine
equalizers in terms of the number of parallel arrows alone, even
for two arrows in $\Delta_{\rm inj}$.
\end{theorem}
\begin{proof}
We first give a uniform lower bound for one finite ordered-set
family. Let $P=[n]$, $Q=[n+1]$, and let
\[
 d_0(j)=j,\qquad d_1(j)=j+1 \quad (1\leq j\leq n).
 \tag{UP.1}
\]
For a finite ordered target $D$, identify an embedding
$P\to D$ with its increasing $n$-tuple of images. An embedding
$g:Q\to D$ gives an edge between its first and last $n$-tuples.
These edges form a simple undirected graph: an edge determines
the increasing $(n+1)$-tuple uniquely. Orient every edge from
the first window to the last window.

This graph has no odd closed walk of length $L<n$. Indeed,
write its successive $n$-tuples as $x^0,\ldots,x^L=x^0$.
For each forward step put $\epsilon_j=1$ and for each backward
step put $\epsilon_j=-1$, and set
$S_0=0$, $S_j=\epsilon_1+\cdots+\epsilon_j$. Since
$\max_j S_j-\min_j S_j\leq L<n$, there is an integer $h$ with
\[
 1+\max_j S_j\leq h\leq n+\min_j S_j.
\]
The coordinate at position $h-S_j$ survives every window shift;
thus $x^j_{h-S_j}=x^0_h$ for every $j$. At the end the distinct
entries of $x^L=x^0$ force $h-S_L=h$, hence $S_L=0$.
An odd number of signs $\pm1$ cannot have sum zero, proving the
claim. Only this unbounded lower bound is needed here.

In fact every closed walk of length $L<n$ has signed forward
minus backward edge count zero, by the same coordinate argument.
Let $R$ be any nonzero unital ring and let
$\mu\in R[\operatorname{Emb}(Q,D)]$ be an affine equalizer.
Regard its support as a finite oriented graph with edge coefficients
in $R$. The equality $\mu d_0=\mu d_1$ says that this chain is a
circulation. Choose a spanning forest of its support. Subtracting
multiples of the fundamental cycles removes each edge outside the
forest, and a circulation supported on a finite forest is zero,
by successively removing its leaves. If the support had fewer
than $n$ edges, each fundamental cycle would have length less
than $n$ and signed augmentation zero. Every circulation on that
support would then have augmentation zero, contrary to
$\operatorname{aug}\mu=1$. Hence
\[
 |\operatorname{supp}\mu|\geq n.
 \tag{UP.2}\label{C4:window-support}
\]
This argument uses only integer signs, so also works for
noncommutative unital rings. For $R=\mathbb F_2$ one can
alternatively note that an odd Eulerian edge set contains an odd
cycle. The same lower bound applies when $P,Q$ carry zero
cocycles in $\mathcal A_r(\mathbb Q)$ and $D$ carries any
cocycle: forgetting labels sends an affine equalizer to one for
(UP.1), preserves augmentation, and can only decrease support.

Take any sequence $n_i\to\infty$ and form the two arrows in
$\mathcal C$ represented by (UP.1) between the zero-cocycle
objects $[n_i]$ and $[n_i+1]$. Suppose they had an affine
equalizer with a finite number $m$ of nonzero terms. Represent
its target and its $m$ arrows by sequences in the factors, keeping their coefficients in $R$ fixed. The
formal equality of the two finite sums of composite arrows is
a finite collection of equalities, after grouping equal terms.
Each of these equalities holds on a member of $\mathcal U$;
their finite intersection is still a member of $\mathcal U$.
On that set of indices the representatives consequently give
an $R$-affine equalizer with at most $m$ terms and
augmentation one. But $n_i>m$ also holds on a member of
$\mathcal U$, contradicting (UP.2).

Every factor has common successors, so the ultraproduct does
too, by choosing coordinatewise successors. The affine criterion
for exactness therefore proves failure of exactness over every
nonzero unital ring. There is also no integral affine equalizer of this
pair, because reduction modulo two would give one. A total
category satisfies the integral pair condition IV, so this
ultraproduct is not total. The identical proof applies to the
ultrapower of $\Delta_{\rm inj}$.
\end{proof}

The preceding theorem concerns the ordinary category
ultraproduct, with all internal objects and arrows and with
ordinary finitely supported module elements. Replacing it by
a selected subcategory or by a construction permitting internal
hyperfinite sums changes the problem and is not covered by an
assertion of modular exactness or totality here.


\subsection{A finite sparse characteristic effect that disappears later}

\begin{proposition}[Proved: targetwise separation need not persist]
\label{C4:sparse-local-global}
In $\Cdelta$, let $V$ be the singleton and let $Q$ have
ordered vertices $a<b<c<x<y$ and edges $\{a,x,y\}$ and
$\{b,c,y\}$. Consider the three strong arrows $V\to Q$
selecting $a,b,c$. There is a strong target $R$ with
$24$ vertices and $14$ edges such that this family has an
affine equalizer over $\mathbb F_2$ in $R$, but none over
$\Q$ in $R$. Nevertheless the same family has an integral
affine equalizer at a larger target in $\Cdelta$.
\end{proposition}
\begin{proof}
Use ten selected vertices in the order
\[
 B<A<C<E<z<D<F<J<H<I
\]
and the following seven increasing triples:
\[
 (z,D,H),\ (A,z,I),\ (B,E,z),\ (A,E,J),\
 (C,D,J),\ (C,F,H),\ (B,F,I).
 \tag{SL.1}
\]
For each triple $(a_i,b_i,c_i)$, add private vertices
$x_i<y_i$, put all private vertices above the ten selected
vertices, and add the edges $a_ix_iy_i,b_ic_iy_i$.
Distinct triples meet in at most one vertex.

Here is a complete finite predimension certificate. For a
subset $S$ of the ten selected vertices, minimizing over
the private vertices of one gadget gives contribution
$-2$ if all three selected vertices of that gadget belong
to $S$, contribution $-1$ if $b_i,c_i\in S$ but $a_i\notin S$,
and contribution zero otherwise. Thus the minimum is
\[
 m(S)=2|S|-2\#\{i:\{a_i,b_i,c_i\}\subseteq S\}
       -\#\{i:\{b_i,c_i\}\subseteq S,\ a_i\notin S\}.
 \tag{SL.2}
\]
The minima of (SL.2) at successive cardinalities
$|S|=0,1,\ldots,10$ are
\[
 0,2,3,4,5,5,6,6,6,6,6.
 \tag{SL.3}
\]
The exact certificate \texttt{code/C4\_sparse\_local\_exception.py}
checks all $1024$ selected subsets and all four private
choices in each gadget; therefore it covers all $2^{24}$
actual subsets. A nonempty subset containing no selected
vertex has no edge and positive predimension. Equations
(SL.2)--(SL.3) prove positivity. Each proposed copy of $Q$
has predimension four. Adding a selected vertex gives
predimension at least five; with its selected triple fixed,
adding other private vertices also strictly raises
predimension, since every other gadget meets the triple
in at most one point. Consequently all seven copies are
strong, and $\delta(R)=6$.

They are all the induced copies of this ordered $Q$ in $R$.
Its two edges must meet in their greatest vertex. Two
edges from different gadgets can meet only at a selected
vertex, which is below their private vertices. They cannot
form such a copy. Within a gadget the increasing order
forces the displayed identification with $Q$.

Let $D$ have one column per copy, recording
$(b_i-a_i,c_i-a_i)$ in two copies of the free module on
the selected vertices. Every vertex other than $z$ occurs
twice in (SL.1), always in the same position; $z$ occurs
once in each position. Hence the sum of the seven columns
is zero modulo two, with augmentation $7=1$ modulo two.
Over $\Q$, a vector $(t_i)_{i=0}^6$ in the kernel must
satisfy $t_i+t_j=0$ whenever two triples share a vertex
other than $z$. The resulting graph of the seven triples
is connected and contains the triangle $0,4,5,0$.
It follows that every $t_i=0$. The exact integer matrix
has Smith factors $1,1,1,1,1,1,2$, as independently
recorded by the certificate. There is therefore no
rational affine equalizer at this target.

We next prove that the targetwise obstruction disappears.
Let $M$ be the strong Fraisse limit and let $\Omega$ be
its strong-singleton orbit. Form the abelian group with
generators $u_t,v_t$ for $t\in\Omega$ and one generator
$h$, with the relations
\[
 u_b-u_a+v_c-v_a=h
 \quad\text{for every strong ordered copy }(a,b,c,x,y)
       \text{ of }Q\text{ in }M.
 \tag{SL.4}
\]
Write $\psi_t=u_t+v_t$. A three-edge star with edges
$axy,a'x'y,bcy$, ordered so that both $a,a'$ are below
$b<c$ and both $x,x'$ lie between $c$ and $y$, has
seven vertices and predimension five. Its two displayed
copies of $Q$ are strong. Subtracting their relations
gives $\psi_a=\psi_{a'}$.

The pair $a,a'$ is not strong in this star; this is why
one star does not suffice to compare an arbitrary strong
pair. Glue two such stars along an outer minimum, obtaining
outer minima $a_0,a_1,a_2$. All other vertices may be put
above these three minima in an order satisfying both
stars. Their free amalgam is positive, and its four
specified copies of $Q$ are strong. Its endpoint pair
$\{a_0,a_2\}$ is strong: if a subset containing this pair
omits $a_1$, its intersections with the two stars each
have predimension at least two, strictly larger when
enlarged beyond the endpoint; if it contains $a_1$, each
intersection contains the two outer minima and has
predimension at least four, so their union has
predimension at least $4+4-2=6>4$. This also has an
independent exhaustive certificate on all $8191$ nonempty
subsets of the thirteen-vertex bridge.

The extension property now gives $\psi_s=\psi_t$ for any
strong pair $s,t$ in $M$. For an arbitrary pair in
$\Omega$, add one singleton freely over a finite strong
set containing it. Each old point and the new point form
a strong pair, so $\psi$ is constant on all of $\Omega$.
Every strong edge $b<c<y$ extends to a strong copy of $Q$
by adding $a<b$ and $c<x<y$. Substituting the constant
$\psi$ into (SL.4) gives
\[
 u_b-u_c=h\qquad\text{for every strong ordered edge }b<c<y.
 \tag{SL.5}
\]
The integral pair affine-equalizer property of the total
category $\Cdelta$, applied to the first two singleton
arrows into an edge, makes a finite integer combination
of (SL.5) have left side zero and right side $h$.
Therefore $h=0$ in the group presented by (SL.4).

Membership in the subgroup generated by the relations
uses only finitely many relations. Writing $h$ as their
integer combination gives a finite integer combination
of copies of $Q$ whose two marginal differences vanish
and whose augmentation is $1$, after changing all signs
if necessary. A finite strong set containing their images
is a target for an integral affine equalizer. Taking a
common successor with $R$ if desired makes it a target
extending the original finite-stage certificate.
\end{proof}

\paragraph{Scope.}
This example separates a matrix calculation at one target
from the existence of an affine equalizer somewhere in the
category. It neither gives a familywise modular-versus-rational
separation in a total category nor changes the status of
the total restriction of C.4.


\subsection{Fixed finite edge covers do not yield all-prime exactness}

\begin{proposition}[Proved: a persistent obstruction for edge covers]
\label{C4:voltage-cover-exclusion}
Let $I$ be a nonempty finite ordered set and let $\sigma,\tau$
be permutations of $I$. For $A\in\Cdelta$, define
$F_{I,\sigma,\tau}(A)$ on the lexicographically ordered set
$A\times I$ by replacing each increasing edge $a<b<c$ by
\[
 \bigl\{(a,i),(b,\sigma(i)),(c,\tau(i))\bigr\},\qquad i\in I.
 \tag{COV.1}
\]
On arrows put $F(f)(v,i)=(f(v),i)$. This is a strong
endofunctor of $\Cdelta$.

There is one fixed three-arrow family in $\Cdelta$ whose
image under every such functor, and under every finite
composition of such functors, has no affine equalizer over
any nonzero unital ring. Consequently a filtered colimit of
copies of $\Cdelta$ with these transition functors remains
nonexact over every nonzero unital ring, although it is total.
\end{proposition}
\begin{proof}
For $U\subseteq F(A)$ write
$m_v=|U\cap(\{v\}\times I)|$ and
$A_t=\{v:m_v\geq t\}$. Above an old edge $a<b<c$, at
most $\min(m_a,m_b,m_c)$ lifted edges lie in $U$, since
all three fiber maps in (COV.1) are injective. Therefore
\[
 \delta(U)\geq \sum_{t=1}^{|I|}\delta(A_t).
 \tag{COV.2}
\]
This proves positivity for every nonempty $U$. If
$B\leq A$ and $U$ properly contains $F(B)$, every $A_t$
contains $B$ and at least one properly contains it.
Strict strongness gives
$\delta(U)>|I|\delta(B)=\delta(F(B))$.
Thus $F(B)\leq F(A)$. Inducedness and functoriality follow
directly from the definition.

The persistent family uses the nine-vertex object $W_*$
with order
\[
 a_0<a_1<a_2<b_0<b_1<b_2<c_0<c_1<c_2
\]
and five edges
\[
 a_0b_0c_0,\quad a_0b_1c_1,\quad a_0b_2c_2,
 \quad a_1b_0c_1,\quad a_2b_1c_2.
 \tag{COV.3}
\]
It has $\delta(W_*)=3$. It is linear, and the minimum
numbers of vertices used by $1,2,3,4,5$ of its edges are,
respectively, $3,5,6,8,9$. Hence every subset with at least
two vertices has predimension at least three. All nine
singletons are strong, and no other proper nonempty subset
is strong. The exact certificate
\texttt{code/C4\_voltage\_cover\_certificate.py}
checks all $512$ subsets and these edge-union bounds.

Let $a_0,a_1,a_2:V\to W_*$ denote the three singleton
arrows with the indicated images. Their marginal-difference
map is injective over every coefficient ring. Indeed, for
a finite collection of $r$ distinct strong copies of $W_*$
in a target, intersections of copies are strong in each
copy and so have at most one vertex. Their five-edge sets
are therefore pairwise disjoint. If $v$ is the number of
vertices appearing in the three selected positions, the
whole union has at most $v+6r$ vertices and at least $5r$
edges. Positivity gives
\[
 0<\delta(\text{union})\leq 2v-3r.
 \tag{COV.4}
\]
If every selected vertex belonged to at least two selected
triples, then $2v\leq3r$, a contradiction. Thus some selected
vertex belongs to exactly one copy. Its coefficient in at
least one of the two marginal differences is $1$ or $-1$.
Applying this argument to the nonzero support of a kernel
element forces a nonzero coefficient to vanish. This proves
injectivity over arbitrary nonzero rings, and excludes an
affine equalizer.

Fix $i\in I$. There is an increasing induced section
$s:W_*\to F(W_*)$ given by
\[
 s(a_j)=(a_j,i),\qquad s(b_j)=(b_j,\sigma(i)),
 \qquad s(c_j)=(c_j,\tau(i)).
 \tag{COV.5}
\]
It preserves all five edges and introduces none. It is
strong: for a subset $U$ properly containing its image,
one has $A_1=W_*$ and some nonempty $A_t$ with $t\geq2$.
Equation (COV.2) and positivity imply
$\delta(U)>\delta(W_*)$.
The singleton section $t:V\to F(V)$ choosing the index
$i$ is strong for the same reason. These sections obey
\[
 F(a_j)t=s a_j\qquad(j=0,1,2).
 \tag{COV.6}
\]
An affine equalizer of the three $F(a_j)$ could therefore
be composed on the right with $t$ and restricted through
$s$, producing an affine equalizer of the original three
arrows with unchanged augmentation. This is impossible
over every nonzero ring.

Composing two covers replaces the ordered fiber by the
lexicographically ordered product and the two permutations
by their componentwise products. Thus every finite
composition has the same form and the same obstruction.
In a filtered colimit of categories, any finite affine
witness, including its finitely many equality checks,
occurs at some finite stage. Applying this to the image
of (COV.3)'s family contradicts the obstruction just proved.

For completeness, filtered colimits preserve totality:
restrict a set action on the colimit category to all stages.
A finite simplicial sphere in a connected component of its
category of elements, together with finitely many paths
connecting its vertices, lifts to a common stage. That
stage's action component is contractible, so the sphere
has a nullhomotopy there and hence in the colimit.
The same finite-diagram argument handles components and
$\pi_1$. Since each stage $\Cdelta$ is total, the claimed
filtered colimit is total.
\end{proof}

\paragraph{Scope.}
The conclusion covers arbitrary finite fiber sizes and
arbitrary fixed permutations in (COV.1), not only cyclic
prime fibers. It does not cover arbitrary endofunctors,
templates that replace edges by larger internal gadgets,
or a modification in which the transition is not such a
finite cover. The total all-prime implication remains open.

